\subsection{CHANGE OF RING}

PROPOSITION 2. Let A, B be two rings, \(\rho: \mathbf{B} \to \mathbf{A}\) a ring homomorphism and E (resp. F) a right (resp. left) A-module. Then there exists one and only one Z-linear mapping

\[
\phi : \rho_ {*} (\mathbf {E}) \otimes_ {\mathbf {B}} \rho_ {*} (\mathbf {F}) \rightarrow \mathbf {E} \otimes_ {\mathbf {A}} \mathbf {F}\tag{6}
\]

such that, for all \(x \in \mathbf{E}\) and \(y \in \mathbf{F}\) , the image under \(\phi\) of the element \(x \otimes y\) of \(\rho_{*}(\mathbf{E}) \otimes_{\mathbf{B}} \rho_{*}(\mathbf{F})\) is the element \(x \otimes y\) of \(\mathbf{E} \otimes_{\mathbf{A}} \mathbf{F}\) ; this \(\mathbf{Z}\) -linear mapping is surjective.

We consider the mapping \((x,y)\mapsto x\otimes y\) of \(\rho_{\star}(E)\times \rho_{*}(F)\) into \(E\otimes_A F\) ; it is Z-bilinear and, for all \(\beta \in B\) , by definition \((x\rho (\beta))\otimes y = x\otimes (\rho (\beta)y)\) , hence conditions (1) of no. 1 hold, whence the existence and uniqueness of \(\phi\) (no. 1, Proposition 1). The latter assertion follows from the fact that the elements \(x\otimes y\) generate the Z-module \(E\otimes_A F\) .

The mapping (6) is called canonical.

COROLLARY. Let \(\mathfrak{a}\) be a two-sided ideal of A such that \(\mathfrak{a}\) is contained in the annihilator of E and in the annihilator of F, so that E (resp. F) has a canonical left (resp. right) \((A/\mathfrak{a})\)-module structure (§ 1, no. 12). Then the canonical homomorphism (6)

\[
\phi : \mathbf {E} \otimes_ {\mathbf {A}} \mathbf {F} \rightarrow \mathbf {E} \otimes_ {\mathbf {A} / \mathfrak {a}} \mathbf {F}
\]

corresponding to the canonical homomorphism \(\rho: \mathbf{A} \to \mathbf{A} / \mathfrak{a}\) is the identity.

For all \(\bar{\alpha} \in \mathbf{A} / \mathfrak{a}\) , all \(x \in \mathbf{E}\) and all \(y \in \mathbf{F}\) , \(x\bar{\alpha} = x\alpha\) (resp. \(\bar{\alpha}y = \alpha y\) ) for all \(\alpha\) such that \(\rho(\alpha) = \bar{\alpha}\) . If \(\mathbf{C} = \mathbf{Z}^{(\mathbb{E} \times \mathbb{F})}\) , the submodule of \(\mathbf{C}\) generated by the elements \((x\alpha, y) - (x, \alpha y)\) is then equal to the submodule generated by the elements \((x\bar{\alpha}, y) - (x, \bar{\alpha}y)\) .

With the hypotheses and notation of Proposition 2, let \(\mathbf{E}'\) be a right B-module, \(\mathbf{F}'\) a left B-module and consider two semi-linear mappings \(u: \mathbf{E}' \to \mathbf{E}\) , \(v: \mathrm{F}' \to \mathrm{F}\) relative to the homomorphism \(\rho: \mathrm{B} \to \mathrm{A}\) ; \(u\) (resp. \(v\) ) can be considered as a B-linear mapping \(\mathrm{E}' \to \rho_*(\mathrm{E})\) (resp. \(\mathrm{F}' \to \rho_*(\mathrm{F})\) ), whence a Z-linear mapping \(w: \mathrm{E}' \otimes_{\mathrm{B}} \mathrm{F}' \to \rho_*(\mathrm{E}) \otimes_{\mathrm{B}} \rho_*(\mathrm{F})\) such that \(w(x' \otimes y') = u(x') \otimes v(y')\) for \(x' \in \mathrm{E}', y' \in \mathrm{F}'\) ; composing the canonical mapping (6) with this mapping, a Z-linear mapping \(w': \mathrm{E}' \otimes_{\mathrm{B}} \mathrm{F}' \to \mathrm{E} \otimes_{\mathrm{A}} \mathrm{F}\) is hence obtained such that \(w'(x' \otimes y') = u(x') \otimes v(y')\) for \(x' \in \mathrm{E}', y' \in \mathrm{F}'\) ; this is the mapping which will normally be denoted by \(u \otimes v\) if no confusion can arise. Clearly \((u, v) \mapsto u \otimes v\) is a Z-bilinear mapping

\[
\operatorname{Hom} _ {\mathbf {B}} \left(\mathrm{E} ^ {\prime}, \rho_ {*} (\mathrm{E})\right) \times \operatorname{Hom} _ {\mathbf {B}} \left(\mathrm{F} ^ {\prime}, \rho_ {*} (\mathrm{F})\right)\rightarrow \operatorname{Hom} _ {\mathbf {Z}} \left(\mathrm{E} ^ {\prime} \otimes_ {\mathbf {B}} \mathrm{F} ^ {\prime}, \mathrm{E} \otimes_ {\mathbf {A}} \mathrm{F}\right).
\]

Moreover, if C is a third ring, \(\sigma: \mathrm{C} \to \mathrm{B}\) a homomorphism, \(\mathrm{E}''\) a right C-module, \(\mathrm{F}''\) a left C-module, \(u': \mathrm{E}'' \to \mathrm{E}'\) and \(v': \mathrm{F}'' \to \mathrm{F}'\) semi-linear mappings relative to \(\sigma\) , then

\[
(u \circ u ^ {\prime}) \otimes (v \circ v ^ {\prime}) = (u \otimes v) \circ (u ^ {\prime} \otimes v ^ {\prime}).
\]

\subsection{OPERATORS ON A TENSOR PRODUCT; TENSOR PRODUCTS AS MULTI-MODULES}

With the hypotheses and notation of no. 1, for every endomorphism \(u\) (resp. \(v\) ) of the A-module \(E\) (resp. \(F\) ), \(u \otimes l_F\) (resp. \(l_E \otimes v\) ) is an endomorphism of the \(Z\) -module \(E \otimes_A F\) ; it follows immediately from (5) (no. 2) that the mapping \(u \mapsto u \otimes l_F\) (resp. \(v \mapsto l_E \otimes v\) ) is a ring homomorphism

\[
\operatorname{End} _ {\mathbf {A}} (\mathbf {E}) \rightarrow \operatorname{End} _ {\mathbf {Z}} (\mathbf {E} \otimes_ {\mathbf {A}} \mathbf {F})
\]

(resp. \(\operatorname{End}_{\mathbf{A}}(\mathbf{F}) \to \operatorname{End}_{\mathbf{Z}}(\mathbf{E} \otimes_{\mathbf{A}} \mathbf{F}))\) ; moreover,

\[
(u \otimes \mathrm{l} _ {\mathrm{F}}) \circ (\mathrm{l} _ {\mathrm{E}} \otimes v) = (\mathrm{l} _ {\mathrm{E}} \otimes v) \circ (u \otimes \mathrm{l} _ {\mathrm{F}}) = u \otimes v\tag{7}
\]

and therefore (§ 1, no. 14) E ⊗\_A F has a canonical left bimodule structure with respect to the rings End\_A(E) and End\_A(F).

This being so, suppose given on E a \(((\mathbf{B}_i^{\prime});\mathbf{A},(\mathbf{C}_j^{\prime}))\) -multimodule structure and on F a \((\mathbf{A},(\mathbf{B}_h^{\prime});(\mathbf{C}_k^{\prime}))\) -multimodule structure (§ 1, no. 14); it amounts to the same to say that ring homomorphisms \(\mathrm{B}_i^\prime \to \mathrm{End}_{\mathbf{A}}(\mathbf{E})\) , \(\mathrm{C}_j^{\prime 0} \to \mathrm{End}_{\mathbf{A}}(\mathbf{E})\) are given with pairwise permutable images, and ring homomorphisms \(\mathrm{B}_h^{\prime \prime} \to \mathrm{End}_{\mathbf{A}}(\mathbf{F})\) , \(\mathrm{C}_k^{\prime \prime 0} \to \mathrm{End}_{\mathbf{A}}(\mathbf{F})\) with pairwise permutable images. If these homomorphisms are composed respectively with the canonical homomorphisms \(\mathrm{End}_{\mathbf{A}}(\mathbf{E}) \to \mathrm{End}_{\mathbf{Z}}(\mathbf{E} \otimes_{\mathbf{A}} \mathbf{F})\) and \(\mathrm{End}_{\mathbf{A}}(\mathbf{F}) \to \mathrm{End}_{\mathbf{Z}}(\mathbf{E} \otimes_{\mathbf{A}} \mathbf{F})\) defined above, it is seen (taking account of (7)) that ring homomorphisms

\[
\mathrm{B} _ {i} ^ {\prime} \rightarrow \operatorname{End} _ {\mathbf {Z}} (\mathrm{E} \otimes_ {\mathrm{A}} \mathrm{F}), \quad \mathrm{C} _ {j} ^ {\prime 0} \rightarrow \operatorname{End} _ {\mathbf {Z}} (\mathrm{E} \otimes_ {\mathrm{A}} \mathrm{F})
\]

\[
\mathrm{B} _ {h} ^ {\prime \prime} \rightarrow \operatorname{End} _ {\mathbf {Z}} (\mathrm{E} \otimes_ {\mathrm{A}} \mathrm{F}), \quad \mathrm{C} _ {k} ^ {\prime \prime 0} \rightarrow \operatorname{End} _ {\mathbf {Z}} (\mathrm{E} \otimes_ {\mathrm{A}} \mathrm{F})
\]

are defined with pairwise permutable images; in other words, there has been defined on \(\mathbf{E} \otimes_{\mathbf{A}} \mathbf{F}\) a \(((\mathbf{B}_i'), (\mathbf{B}_h''); (\mathbf{C}_j'), (\mathbf{C}_k'')\) -multimodule structure; it is this multimodule which is also called the tensor product (relative to A) of the \(((\mathbf{B}_i^{\prime});\mathbf{A},(\mathbf{C}_j^{\prime}))\) -multimodule E and the \((\mathbf{A},(\mathbf{B}_h^{\prime\prime});(\mathbf{C}_k^{\prime\prime}))\) -multimodule F. This multi-module is the solution of a universal mapping problem analogous to that considered in no. 1; to be precise:

PROPOSITION 3. Let G be a \(((\mathbf{B}_i'), (\mathbf{B}_h''); (\mathbf{C}_j'), (\mathbf{C}_k''))-multimodule.\)

\begin{enumerate}
\def\labelenumi{(\alph{enumi})}
\tightlist
\item
  Let \(g\) be a linear mapping of the multimodule \(\mathbf{E} \otimes_{\mathbf{A}} \mathbf{F}\) into \(\mathbf{G}\) . The mapping \(f: (x, y) \mapsto g(x \otimes y)\) of \(\mathbf{E} \times \mathbf{F}\) into \(\mathbf{G}\) is \(\mathbf{Z}\) -bilinear and satisfies relations (1) of no. 1 and the conditions
\end{enumerate}

\[
\left\{ \begin{array}{l l} f (\mu_ {i} ^ {\prime} x, y) = \mu_ {i} ^ {\prime} f (x, y), & f (x \nu_ {j} ^ {\prime}, y) = f (x, y) \nu_ {j} ^ {\prime} \\ f (x, \mu_ {h} ^ {\prime \prime} y) = \mu_ {h} ^ {\prime \prime} f (x, y), & f (x, y \nu_ {k} ^ {\prime \prime}) = f (x, y) \nu_ {k} ^ {\prime \prime} \end{array} \right.\tag{8}
\]

for all \(x \in \mathbf{E}, y \in \mathbf{F}\) , \(\mu_i' \in \mathbf{B}_i'\) , \(\nu_j' \in \mathbf{C}_j'\) , \(\mu_h'' \in \mathbf{B}_h''\) , \(\nu_k'' \in \mathbf{C}_{k}''\) , \(i, j, h, k\) arbitrary.

\begin{enumerate}
\def\labelenumi{(\alph{enumi})}
\setcounter{enumi}{1}
\tightlist
\item
  Conversely, let \(f\) be a \(\mathbf{Z}\) -bilinear mapping of \(\mathbf{E} \times \mathbf{F}\) into \(\mathbf{G}\) satisfying conditions (1) (no. 1) and (8). Then there exists one and only one linear mapping \(g\) of the multimodule \(\mathbf{E} \otimes_{\mathbf{A}} \mathbf{F}\) into the multimodule \(\mathbf{G}\) such that \(f(x, y) = g(x \otimes y)\) for \(x \in \mathbf{E}, y \in \mathbf{F}\) .
\end{enumerate}

Assertion (a) follows immediately from the definition of the multimodule structure on \(\mathbf{E} \otimes_{\mathbf{A}} \mathbf{F}\) , since for example \((x \otimes y)\nu_j' = (x\nu_j') \otimes y\) . To prove (b), we note first that Proposition 1 of no. 1 gives the existence and uniqueness of a \(\mathbf{Z}\) -linear mapping \(g\) such that \(g(x \otimes y) = f(x, y)\) for \(x \in \mathbf{E}\) , \(y \in \mathbf{F}\) ; all that is needed is to verify that \(g\) is linear for the multimodule structures. As the elements \(x \otimes y\) generate the \(\mathbf{Z}\) -module \(\mathbf{E} \otimes_{\mathbf{A}} \mathbf{F}\) , it suffices to verify the relations \(g(\mu'(x \otimes y)) = \mu_i' g(x \otimes y)\) and their analogues; but this follows immediately from the formula \(g(x \otimes y) = f(x, y)\) and relations (8).

SCHOLIUM. An element of \(\mathbf{E} \otimes_{\mathbf{A}} \mathbf{F}\) may in general be written in several ways in the form \(\sum_{i} (x_i \otimes y_i)\) , where \(x_i \in \mathbf{E}\) and \(y_i \in \mathbf{F}\) ; but to define a linear mapping \(g\) of the multimodule \(\mathbf{E} \otimes_{\mathbf{A}} \mathbf{F}\) into a multimodule \(\mathbf{G}\) , there is no need to verify that, if \(\sum_{i} (x_i \otimes y_i) = \sum_{j} (x_j' \otimes y_j')\) , then \(\sum_{i} g(x_i \otimes y_i) = \sum_{j} g(x_j' \otimes y_j')\) ; it suffices to be given \(g(x \otimes y)\) for \(x \in \mathbf{E}\) and \(y \in \mathbf{F}\) and to verify that \((x, y) \mapsto g(x \otimes y)\) is \(\mathbf{Z}\) -bilinear and satisfies conditions (1) (no. 1) and (8).

Let \(\mathbf{E}'\) be a \(((\mathbf{B}_i'),\mathbf{A},(\mathbf{C}_j')\) -multimodule, \(\mathbf{F}'\) an \((\mathbf{A},(\mathbf{B}_h'');(\mathbf{C}_k'')\) -multimodule and \(u:\mathbf{E}\to \mathbf{E}',v:\mathbf{F}\to \mathbf{F}'\) linear mappings of multimodules; it follows immediately from the definitions (no. 2) that \(u\otimes v\) is a linear mapping of the multimodule \(\mathbf{E}\otimes_{\mathbf{A}}\mathbf{F}\) into the multimodule \(\mathbf{E}'\otimes_{\mathbf{A}}\mathbf{F}'\) .

With E always denoting a right A-module, let \(_{s}A_{d}\) denote the ring A considered as an (A, A)-bimodule ( \(\S\) 1, no. 14, Example 1); by the above, the tensor product \(\mathbf{E} \otimes_{\mathbf{A}} (\mathbf{s}\mathbf{A}_{d})\) has a canonical right A-module structure such that \((x \otimes \lambda)\mu = x \otimes (\lambda\mu)\) for \(x \in E\) , \(\lambda \in A\) , \(\mu \in A\) . The mapping \((x, \lambda) \mapsto x\lambda\) of \(\mathbf{E} \times (\mathbf{s}\mathbf{A}_{d})\) into E is Z-bilinear and satisfies conditions (1) (no. 1) and (8)

(where, in the latter, the \(\mathbf{B}_i'\) , \(\mathbf{C}_j'\) and \(\mathbf{B}_h''\) are absent and the family \((\mathbf{C}_k'')\) reduces to \(\mathbf{A}\) ); hence (Proposition 3), there exists an A-linear mapping \(g\) (called canonical) of \(\mathbf{E} \otimes_{\mathbf{A}} (s\mathbf{A}_d)\) into \(\mathbf{E}\) such that \(g(x \otimes \lambda) = x\lambda\) for \(x \in \mathbf{E}\) , \(\lambda \in \mathbf{A}\) .

PROPOSITION 4. If E is a right A-module, the mapping \(h: x \mapsto x \otimes 1\) of E into \(\mathbf{E} \otimes_{\mathbf{A}} (s\mathbf{A}_d)\) is a right A-module isomorphism, whose inverse isomorphism \(g\) is such that \(g(x \otimes \lambda) = x\lambda\) for \(x \in \mathbf{E}, \lambda \in \mathbf{A}\) .

If \(g\) is the canonical mapping, \(g \circ h\) is the identity mapping \(1_{\mathbf{E}}\) and \(h \circ g\) coincides with the identity mapping of \(\mathbf{E} \otimes_{\mathbf{A}} (s\mathbf{A}_d)\) onto itself for elements of the form \(x \otimes y\) , which generate the latter \(\mathbf{Z}\) -module; hence the conclusion.

We shall normally write \(E \otimes_{A} A\) instead of \(E \otimes_{A} (s A_{d})\) and shall often identify \(E \otimes_{A} A\) with E by means of the above canonical isomorphisms. Observe that, if E also has a (left or right) B-module structure which is compatible with its right A-module structure, g and h are also isomorphisms for the B-module structures on E and \(E \otimes_{A} A\) (and hence multimodule isomorphisms).

Now let F be a left A-module; \((sA_{d}) \otimes_{A} F\) (also denoted by A \(\otimes_{A} F\) ) then has a canonical left A-module structure and as in Proposition 4 a canonical isomorphism is defined of A \(\otimes_{A} F\) onto F mapping \(\lambda \otimes x\) to \(\lambda x\) , and its inverse isomorphism is \(x \mapsto 1 \otimes x\) .

In particular there exists a canonical isomorphism of the (A, A)-bimodule \((_{s}\mathbf{A}_{d})\otimes_{\mathbf{A}}(_{s}\mathbf{A}_{d})\) onto \({}_{s}\mathbf{A}_{d}\) which maps \(\lambda \otimes \mu\) to \(\lambda \mu\) .

\subsection{TENSOR PRODUCT OF TWO MODULES OVER A COMMUTATIVE RING}

Let C be a commutative ring; for every C-module E, the module structure on E is compatible with itself (§ 1, no. 14). If E and F are two C-modules, the considerations of no. 4 then allow us to define two C-module structures on the tensor product E ⊗c F, respectively such that γ(x ⊗ y) = (γx) ⊗ y and such that γ(x ⊗ y) = x ⊗ (γy); but as, by Definition 1 of no. 1, in this case (γx) ⊗ y = x ⊗ (γy), these two structures are the same. Henceforth when we speak of E ⊗c F as a C-module, we shall mean with the structure thus defined, unless otherwise stated. The canonical isomorphism.

\[
\sigma : \mathbf {F} \otimes_ {\mathbf {c}} \mathbf {E} \rightarrow \mathbf {E} \otimes_ {\mathbf {c}} \mathbf {F}
\]

(no. 1, Corollary 2 to Proposition 1) is then a C-module isomorphism.

It follows from this definition that, if \((a_{\lambda})_{\lambda\in\mathbb{L}}\) (resp. \((b_{\mu})_{\mu\in\mathbb{M}}\) ) is a generating system of the C-module E (resp. F), \((a_{\lambda}\otimes b_{\mu})\) is a generating system of the C-module \(E\otimes_{C}F\) ; in particular, if E and F are finitely generated C-modules, so is \(E\otimes_{C}F\) .

For every C-module G, the Z-bilinear mappings \(f\) of \(\mathbf{E} \times \mathbf{F}\) into \(\mathbf{G}\) for which

\[
f (\gamma x, y) = f (x, \gamma y) = \gamma f (x, y) \quad \text { for } x \in \mathbf {E}, y \in \mathbf {F}, \gamma \in \mathbf {C}\tag{9}
\]

are then called C-bilinear and form a C-module denoted by \(\mathcal{L}_2(\mathrm{E},\mathrm{F};\mathrm{G})\) ; Proposition 3 (no. 4) defines a canonical C-module isomorphism (cf.~§1, no. 14, Remark 1).

\[
\mathcal {L} _ {2} (\mathrm{E}, \mathrm{F}; \mathrm{G}) \rightarrow \operatorname{Hom} _ {\mathbf {c}} (\mathrm{E} \otimes_ {\mathbf {c}} \mathrm{F}, \mathrm{G}).\tag{10}
\]

Let \(E'\) , \(F'\) be two C-modules and \(u: E \to E'\) , \(v: F \to F'\) two C-linear mappings; then (no. 4) \(u \otimes v\) is a C-linear mapping of \(E \otimes_{\mathbb{C}} F\) into \(E' \otimes_{\mathbb{C}} F'\) . Further, it is immediate that \((u, v) \mapsto u \otimes v\) is a C-bilinear mapping of \(\operatorname{Hom}_{\mathbb{C}}(\mathbf{E}, \mathbf{E}') \times \operatorname{Hom}_{\mathbb{C}}(\mathbf{F}, \mathbf{F}')\) into \(\operatorname{Hom}_{\mathbb{C}}(\mathbf{E} \otimes_{\mathbb{C}} \mathbf{F}, \mathbf{E}' \otimes_{\mathbb{C}} \mathbf{F}')\) ; hence there canonically corresponds to it a C-linear mapping, called canonical:

\[
\operatorname{Hom} _ {\mathbf {c}} (\mathbf {E}, \mathbf {E} ^ {\prime}) \otimes_ {\mathbf {c}} \operatorname{Hom} _ {\mathbf {c}} (\mathbf {F}, \mathbf {F} ^ {\prime}) \to \operatorname{Hom} _ {\mathbf {c}} (\mathbf {E} \otimes_ {\mathbf {c}} \mathbf {F}, \mathbf {E} ^ {\prime} \otimes_ {\mathbf {c}} \mathbf {F} ^ {\prime})\tag{11}
\]

which associates with every element \(u \otimes v\) of the tensor product

\[
\operatorname{Hom} _ {\mathbf {c}} (\mathbf {E}, \mathbf {E} ^ {\prime}) \otimes_ {\mathbf {c}} \operatorname{Hom} _ {\mathbf {c}} (\mathbf {F}, \mathbf {F} ^ {\prime})
\]

the linear mapping \(u \otimes v\) . Note that the canonical mapping (11) is not necessarily injective nor surjective (§ 4, Exercise 2).

Remarks. (1) Let A, B be two commutative rings, \(\rho:B\to A\) a ring homomorphism and E and F two A-modules; then the canonical mapping (6) of no. 3 is a B-linear mapping

\[
\rho_ {*} (\mathbf {E}) \otimes \rho_ {*} (\mathbf {F}) \rightarrow \rho_ {*} (\mathbf {E} \otimes_ {\mathbf {A}} \mathbf {F}).\tag{12}
\]

\begin{enumerate}
\def\labelenumi{(\arabic{enumi})}
\setcounter{enumi}{1}
\tightlist
\item
  What has been said in this no. can be generalized to the following case: let E be a right A-module, F a left A-module, C a commutative ring and \(\rho:C\to A\) a homomorphism of C into A such that \(\rho(\mathbf{C})\) is contained in the centre of A (cf.~III, §1, no. 3). We may then consider the C-modules \(\rho_{*}(\mathbf{E})\) and \(\rho_{*}(\mathbf{F})\) and the hypothesis on \(\rho\) implies that these C-module structures are compatible respectively with the A-module structures on E and F (§1, no. 14). The tensor product \(E\otimes_{A}F\) is thus (by virtue of no. 4) given two C-module structures such that \(\gamma(x\otimes y)=(x\rho(\gamma))\otimes y\) and
\end{enumerate}

\[
\gamma (x \otimes y) = x \otimes (\rho (\gamma) y)
\]

respectively for \(\gamma\in\mathbf{C}\) , \(x\in\mathbf{E}\) , \(y\in\mathbf{F}\) and Definition 1 (no. 1) shows also that these two structures are identical. If \(\mathbf{E}'\) (resp. \(\mathbf{F}'\) ) is a right (resp. left) A-module and \(u:\mathbf{E}\to\mathbf{E}'\) , \(v:\mathbf{F}\to\mathbf{F}'\) are two A-linear mappings, then \(u\otimes v:\mathbf{E}\otimes_{\mathbf{A}}\mathbf{F}\to\mathbf{E}'\otimes_{\mathbf{A}}\mathbf{F}'\) is C-linear for the C-module structures just defined; the mapping \((u,v)\mapsto u\otimes v\) :

\[
\operatorname{Hom} _ {\mathbf {A}} (\mathrm{E}, \mathrm{E} ^ {\prime}) \times \operatorname{Hom} _ {\mathbf {A}} (\mathrm{F}, \mathrm{F} ^ {\prime}) \rightarrow \operatorname{Hom} _ {\mathbf {C}} (\mathrm{E} \otimes_ {\mathbf {A}} \mathrm{F}, \mathrm{E} ^ {\prime} \otimes_ {\mathbf {A}} \mathrm{F} ^ {\prime})
\]

is C-bilinear (for the C-module structures on \(\mathrm{Hom}_{\mathbf{A}}(\mathbf{E},\mathbf{E}^{\prime})\) and \(\mathrm{Hom}_{\mathbf{A}}(\mathbf{F},\mathbf{F}^{\prime})\)

defined in § 1, no. 14, Remark 1), whence we also derive a C-linear mapping, called canonical

\[
\operatorname{Hom} _ {\mathbf {A}} (\mathbf {E}, \mathbf {E} ^ {\prime}) \otimes_ {\mathbf {C}} \operatorname{Hom} _ {\mathbf {A}} (\mathbf {F}, \mathbf {F} ^ {\prime}) \rightarrow \operatorname{Hom} _ {\mathbf {C}} (\mathbf {E} \otimes_ {\mathbf {A}} \mathbf {F}, \mathbf {E} ^ {\prime} \otimes_ {\mathbf {A}} \mathbf {F} ^ {\prime}).\tag{13}
\]

\begin{enumerate}
\def\labelenumi{(\arabic{enumi})}
\setcounter{enumi}{2}
\tightlist
\item
  Let \(\mathbf{A}\) be an integral domain and \(\mathbf{K}\) its field of fractions. If \(\mathbf{E}\) and \(\mathbf{F}\) are two vector \(\mathbf{K}\) -spaces, the canonical mapping
\end{enumerate}

\[
\left(\mathbf {E} _ {[ \mathbf {A} ]}\right) \otimes_ {\mathbf {A}} \left(\mathbf {F} _ {[ \mathbf {A} ]}\right)\rightarrow \mathbf {E} \otimes_ {\mathbf {K}} \mathbf {F}
\]

(no. 3 and §1, no. 13) is bijective. It suffices (no. 4) to prove that if \(f\) is an A-bilinear mapping of \(\mathbf{E} \times \mathbf{F}\) into a vector K-space \(\mathbf{G}\) , \(f\) is also K-bilinear. Now, for all \(\alpha \neq 0\) in \(\mathbf{A}\) .

\[
\alpha f (\alpha^ {- 1} x, y) = f (x, y) = \alpha f (x, \alpha^ {- 1} y)
\]

whence

\[
f (\alpha^ {- 1} x, y) = f (x, \alpha^ {- 1} y) = \alpha^ {- 1} f (x, y)
\]

since G is a vector K-space.

\subsection{PROPERTIES OF E ⊗\_A F RELATIVE TO EXACT SEQUENCES}

PROPOSITION 5. Let E, E', E'' be right A-modules, F a left A-module and

\[
\mathrm{E} ^ {\prime} \stackrel {u} {\longrightarrow} \mathrm{E} \stackrel {v} {\longrightarrow} \mathrm{E} ^ {\prime \prime} \longrightarrow 0\tag{14}
\]

an exact sequence of linear mappings. Writing \(\bar{u} = u \otimes 1_{\mathbb{F}}, \bar{v} = v \otimes 1_{\mathbb{F}}\) , the sequence

\[
\mathrm{E} ^ {\prime} \otimes_ {\mathrm{A}} \mathrm{F} \xrightarrow {\bar {u}} \mathrm{E} \otimes_ {\mathrm{A}} \mathrm{F} \xrightarrow {\bar {v}} \mathrm{E} ^ {\prime \prime} \otimes_ {\mathrm{A}} \mathrm{F} \longrightarrow 0\tag{15}
\]

of Z-homomorphisms is exact.

By virtue of no. 2, formula (5), \(\bar{v} \circ \bar{u} = (v \circ u) \otimes 1_{\mathbf{F}} = 0\) ; the image \(\mathrm{H} = \bar{u} (\mathbf{E}' \otimes \mathbf{F})\) is contained in the kernel \(\mathbf{L} = \operatorname{Ker}(\bar{v})\) ; by passing to the quotient, we therefore derive from \(\bar{v}\) a \(\mathbf{Z}\) -linear mapping \(f\) of the cokernel \(\mathbf{M} = (\mathbf{E} \otimes \mathbf{F}) / \mathrm{H}\) of \(\bar{u}\) into \(\mathbf{E}'' \otimes \mathbf{F}\) ; it must be proved that \(f\) is bijective and it will hence suffice to define a \(\mathbf{Z}\) -linear mapping \(g: \mathbf{E}'' \otimes \mathbf{F} \to \mathbf{M}\) such that \(g \circ f\) and \(f \circ g\) are the identity mappings.

Let \(x'' \in E'', y \in F\) ; by hypothesis there exists \(x \in E\) such that \(v(x) = x''\) . We show that, if \(x_1, x_2\) are two elements of \(E\) such that \(v(x_1) = v(x_2) = x''\) and \(\phi: E \otimes F \to M\) is the canonical mapping, then \(\phi(x_1 \otimes y) = \phi(x_2 \otimes y)\) . It suffices to prove that if \(v(x) = 0\) then \(\phi(x \otimes y) = 0\) , which follows from the fact that \(x = u(x')\) with \(x' \in E'\) , whence \(x \otimes y = u(x') \otimes y = \bar{u}(x' \otimes y) \in H\) . If \((x'', y)\) is mapped to the unique value of \(\phi(x \otimes y)\) for all \(x \in E\) such that \(v(x) = x''\) , a mapping is defined of \(E'' \times F\) into \(M\) ; this mapping is Z-bilinear and satisfies conditions (1) (no. 1), since \(v(x\lambda) = x''\lambda\) and \((x\lambda) \otimes y = x \otimes (\lambda y)\) for \(x \in E\) ; hence there is a Z-linear mapping \(g\) of \(E'' \otimes F\) into \(M\) such that \(g(x'' \otimes y) = \phi(x \otimes y)\) for \(y \in F, x \in E\) and \(x'' = v(x)\) . This definition further proves that \(f \circ g\) coincides with the identity mapping for the elements of

\(E'' \otimes F\) of the form \(x'' \otimes y\) and hence \(f \circ g\) is the identity mapping of \(E'' \otimes F\) ; on the other hand, for \(x \in E\) and \(y \in F\) , \(f(\phi(x \otimes y)) = v(x) \otimes y\) by definition, hence \(g(f(\phi(x \otimes y))) = \phi(x \otimes y)\) and, as the elements of the form \(\phi(x \otimes y)\) generate M, \(g \circ f\) is the identity mapping of M.

COROLLARY. Let F, F', F'' be left A-modules, E a right A-module and

\[
\mathrm{F} ^ {\prime} \stackrel {s} {\longrightarrow} \mathrm{F} \stackrel {t} {\longrightarrow} \mathrm{F} ^ {\prime \prime} \longrightarrow 0\tag{16}
\]

an exact sequence of linear mappings. Writing \(\bar{s} = 1_{\mathbb{E}} \otimes s\) , \(\bar{t} = 1_{\mathbb{E}} \otimes t\) , the sequence of \(\mathbf{Z}\) -homomorphisms

\[
\mathrm{E} \otimes_ {\mathrm{A}} \mathrm{F} ^ {\prime} \stackrel {\hat {s}} {\longrightarrow} \mathrm{E} \otimes_ {\mathrm{A}} \mathrm{F} \stackrel {\hat {t}} {\longrightarrow} \mathrm{E} \otimes_ {\mathrm{A}} \mathrm{F} ^ {\prime \prime} \longrightarrow 0\tag{17}
\]

is exact.

When E (resp. F) is considered as a left (resp. right) A \(^{0}\) -module, F ⊗A \(^{0}\) E is identified with E ⊗A F and there are analogous identifications for F' ⊗A \(^{0}\) E and F'\,' ⊗A \(^{0}\) E (no. 1, Corollary 2 to Proposition 1); the corollary then follows immediately from Proposition 5.

Remark. Note that in general, if \(E'\) is a submodule of a right A-module E and \(j:E'\to E\) the canonical injection, the canonical mapping

\[
j \otimes 1 _ {\mathbf {F}}: \mathbf {E} ^ {\prime} \otimes \mathbf {F} \rightarrow \mathbf {E} \otimes \mathbf {F}
\]

is not necessarily injective. In other words, for an exact sequence

\[
0 \longrightarrow \mathrm{E} ^ {\prime} \stackrel {u} {\longrightarrow} \mathrm{E} \stackrel {v} {\longrightarrow} \mathrm{E} ^ {\prime \prime} \longrightarrow 0\tag{18}
\]

it cannot in general be concluded that the sequence

\[
0 \longrightarrow \mathrm{E} ^ {\prime} \otimes \mathrm{F} \stackrel {u} {\longrightarrow} \mathrm{E} \otimes \mathrm{F} \stackrel {v} {\longrightarrow} \mathrm{E} ^ {\prime \prime} \otimes \mathrm{F} \longrightarrow 0\tag{19}
\]

is exact.

Take for example A = Z, E = Z, \(E' = 2Z\) , F = Z/2Z. As \(E'\) is isomorphic to E, \(E' \otimes F\) is isomorphic to E \(\otimes\) F which is itself isomorphic to F (no. 4, Proposition 4). But for all \(x' = 2x \in E'\) (where \(x \in E\) ) and all \(y \in F\) , \(j(x') \otimes y = (2x) \otimes y = x \otimes (2y) = 0\) , since 2y = 0, and the canonical image of \(E' \otimes F\) in E \(\otimes\) F reduces to 0.

In other words, care must be taken to distinguish, for a submodule \(E'\) of E and an element \(x \in E'\) , between the element \(x \otimes y\) ``calculated in \(E' \otimes F\)'' and the element \(x \otimes y\) ``calculated in \(E \otimes F\)'' (in other words, the element \(j(x) \otimes y\) ).

We shall study later, under the name of flat modules, the modules F such that the sequence (19) is exact for every exact sequence (18) (Commutative Algebra, I, § 2).

PROPOSITION 6. Given two exact sequences (14) and (16), the homomorphism \(v \otimes t: \mathbf{E} \otimes_{\mathbf{A}} \mathbf{F} \to \mathbf{E}'' \otimes_{\mathbf{A}} \mathbf{F}''\) is surjective and its kernel is equal to

\[
\operatorname{Im} (u \otimes \mathrm{l} _ {\mathrm{F}}) + \operatorname{Im} (\mathrm{l} _ {\mathrm{E}} \otimes s)
\]

Now \(v \otimes t = (v \otimes 1_{\mathbb{F}''}) \circ (1_{\mathbb{E}} \otimes t)\) (no. 2, formula (5)) and \(v \otimes t\) is therefore surjective, being the composition of two surjective homomorphisms by virtue of Proposition 5 and its Corollary. On the other hand, for \(z \in \mathbb{E} \otimes \mathbb{F}\) to be in the kernel of \(v \otimes t\) , it is necessary and sufficient that \((1_{\mathbb{E}} \otimes t) (z)\) belong to the kernel of \(v \otimes 1_{\mathbb{F}''}\) , that is, by virtue of (15) to the image of

\[
u \otimes 1 _ {F ^ {\prime \prime}}: \mathrm{E} ^ {\prime} \otimes \mathrm{F} ^ {\prime \prime} \rightarrow \mathrm{E} \otimes \mathrm{F} ^ {\prime \prime}.
\]

But as the homomorphism \(t: \mathbf{F} \to \mathbf{F}''\) is surjective, so is

\[
\mathbf {l} _ {\mathbf {E} ^ {\prime}} \otimes t: \mathbf {E} ^ {\prime} \otimes \mathbf {F} \rightarrow \mathbf {E} ^ {\prime} \otimes \mathbf {F} ^ {\prime \prime}
\]

by the Corollary to Proposition 5, hence the condition on \(z\) reduces to the existence of an \(a \in \mathbf{E}' \otimes \mathbf{F}\) such that

\[
(1 _ {\mathrm{E}} \otimes t) (z) = (u \otimes t) (a).
\]

Let \(b = z - (u \otimes 1_{\mathbb{F}})(a)\) ; then \((1_{\mathbb{E}} \otimes t)(b) = 0\) , and by virtue of (17), \(b\) belongs to the image of \(1_{\mathbb{E}} \otimes s\) , which proves the proposition.

In other words:

COROLLARY 1. Let \(\mathbf{E}'\) be a submodule of a right A-module \(\mathbf{E}\) , \(\mathbf{F}'\) a submodule of a left A-module \(\mathbf{F}\) and \(\operatorname{Im}(\mathbf{E}' \otimes_{\mathbf{A}} \mathbf{F})\) and \(\operatorname{Im}(\mathbf{E} \otimes_{\mathbf{A}} \mathbf{F}')\) the sub- \(\mathbf{Z}\) -modules of \(\mathbf{E} \otimes_{\mathbf{A}} \mathbf{F}\) , the respective images of the canonical mappings \(\mathbf{E}' \otimes_{\mathbf{A}} \mathbf{F} \to \mathbf{E} \otimes_{\mathbf{A}} \mathbf{F}\) ,

\[
\mathbf {E} \otimes_ {\mathbf {A}} \mathbf {F} ^ {\prime} \rightarrow \mathbf {E} \otimes_ {\mathbf {A}} \mathbf {F}.
\]

Then there is a canonical Z-module isomorphism

\[
\pi : (\mathrm{E} / \mathrm{E} ^ {\prime}) \otimes_ {\mathrm{A}} (\mathrm{F} / \mathrm{F} ^ {\prime}) \rightarrow (\mathrm{E} \otimes_ {\mathrm{A}} \mathrm{F}) / (\operatorname{Im} (\mathrm{E} ^ {\prime} \otimes_ {\mathrm{A}} \mathrm{F}) + \operatorname{Im} (\mathrm{E} \otimes_ {\mathrm{A}} \mathrm{F} ^ {\prime}))\tag{20}
\]

such that, for \(\xi \in \mathbf{E} / \mathbf{E}'\) , \(\eta \in \mathbf{F}\) , \(\pi (\xi \otimes \eta)\) is the class of all elements \(x \otimes y \in \mathbf{E} \otimes_{\mathbf{A}} \mathbf{F}\) such that \(x \in \xi\) and \(y \in \eta\) .

Note that when E is a \(((B_{i}^{\prime}); A, (C_{j}^{\prime}))\) -multimodule, F a \((A, (B_{h}^{\prime\prime}); (C_{k}^{\prime\prime}))\) -multimodule and \(E^{\prime}\) and \(F^{\prime}\) submultimodules of E and F respectively, the isomorphism (20) is an isomorphism for the \(((B_{i}^{\prime}), (B_{h}^{\prime\prime}); (C_{j}^{\prime}), (C_{k}^{\prime\prime}))\) -multimodule structures of the two sides (no. 3).

COROLLARY 2. Let \(\mathfrak{a}\) be a right ideal of \(\mathbf{A}\) , \(\mathbf{F}\) a left A-module and \(\mathfrak{a}\mathbf{F}\) the sub- \(\mathbf{Z}\) -module of \(\mathbf{F}\) generated by the elements of the form \(\lambda x\) , where \(\lambda \in \mathfrak{a}\) and \(x \in \mathbf{F}\) . Then there is a canonical \(\mathbf{Z}\) -module isomorphism

\[
\pi : (\mathrm{A} / \mathfrak {a}) \otimes_ {\mathrm{A}} \mathrm{F} \rightarrow \mathrm{F} / \mathfrak {a} \mathrm{F}\tag{21}
\]

such that, for all \(\bar{\lambda} \in \mathbf{A} / \mathfrak{a}\) and all \(x \in \mathbf{F}\) , \(\pi(\bar{\lambda} \otimes x)\) is the class mod. \(\mathfrak{aF}\) of \(\lambda x\) , where \(\lambda \in \bar{\lambda}\) .

In particular, for A = Z, it is seen that for every integer n and every Z-module F, \((\mathbf{Z}/n\mathbf{Z}) \otimes_{\mathbf{Z}} F\) is canonically identified with the quotient Z-module F/nF.

COROLLARY 3. Let A be a commutative ring, a an ideal of A and E and F two A-modules such that a is contained in the annihilator of F. Then the (A/a)-modules E ⊗\_A F and (E/aE) ⊗\_\{A/a\} F are canonically isomorphic.

F and \(E \otimes_{A} F\) are annihilated by a and hence have canonical (A/a)-module structures (§1, no. 12) and if we write \(E' = aE\) , then \(\operatorname{Im}(E' \otimes_{A} F) = 0\) ; then there is a canonical isomorphism (20) of \(E \otimes_{A} F\) onto \((\mathrm{E}/a\mathrm{E}) \otimes_{A} F\) and the latter is itself identical with \((\mathrm{E}/a\mathrm{E}) \otimes_{A/a} F\) (no. 3, Corollary to Proposition 2).

COROLLARY 4. Let \(\mathfrak{a}, \mathfrak{b}\) be two ideals in a commutative ring \(\mathbf{C}\) ; the \(\mathbf{C}\) -module \((\mathbf{C} / \mathfrak{a}) \otimes_{\mathbb{C}} (\mathbf{C} / \mathfrak{b})\) is then canonically isomorphic to \(\mathbf{C}/(\mathfrak{a} + \mathfrak{b})\) .

\subsection{TENSOR PRODUCTS OF PRODUCTS AND DIRECT SUMS}

Let \((\mathbf{E}_{\lambda})_{\lambda \in L}\) be a family of right A-modules, \((\mathbf{F}_{\mu})_{\mu \in M}\) a family of left A-modules and consider the product modules \(C = \prod_{\lambda \in L} E_{\lambda}\) , \(D = \prod_{\mu \in M} F_{\mu}\) . The mapping \(((x_{\lambda}), (y_{\mu})) \mapsto (x_{\lambda} \otimes y_{\mu})\) of \(C \times D\) into the product Z-module

\[
\prod_ {(\lambda , \mu) \in L \times M} (E _ {\lambda} \otimes_ {A} F _ {\mu})
\]

is Z-bilinear and obviously satisfies conditions (1) (no. 1). Thus there exists (no. 1, Proposition 1) a Z-linear mapping, called canonical

\[
f: \left(\prod_ {\lambda \in L} E _ {\lambda}\right) \otimes_ {A} \left(\prod_ {\mu \in M} F _ {\mu}\right)\rightarrow \prod_ {(\lambda , \mu) \in L \times M} (E _ {\lambda} \otimes_ {A} F _ {\mu})\tag{22}
\]

such that \(f((x_{\lambda}) \otimes (y_{\mu})) = (x_{\lambda} \otimes y_{\mu})\) .

When \(C = R^{L}\) , \(D = S^{M}\) , R (resp. S) being a right (resp. left) A-module, the canonical mapping (22) associates with every tensor product \(u \otimes v\) , where u is a mapping of L into R and v a mapping of M into S, the mapping \((\lambda, \mu) \mapsto u(\lambda) \otimes v(\mu)\) of \(L \times M\) into \(R \otimes_{A} S\) ; even in this case the canonical mapping (22) is in general neither injective nor surjective (Exercise 3; cf.~Corollary 3 to Proposition 7).

When the \(E_{\lambda}\) are \(((B_{i}^{\prime}); A, (C_{j}^{\prime}))\) -multimodules and the \(F_{\mu}(A, (B_{h}^{\prime\prime}); (C_{k}^{\prime\prime}))\) -multimodules, the homomorphism (22) is also a homomorphism for the \(((B_{i}^{\prime}), (B_{h}^{\prime\prime}); (C_{j}^{\prime}), (C_{k}^{\prime\prime}))\) -multimodule structures of the two sides.

Consider now the submodule \(\mathbf{E} = \bigoplus_{\lambda \in L} \mathbf{E}_{\lambda}\) (resp. \(\bigoplus_{\mu \in M} F_{\mu}\) ) of C (resp. D); the canonical injections \(\mathbf{E} \to \mathbf{C}\) , \(\mathbf{F} \to \mathbf{D}\) define canonically a Z-linear mapping \(\mathbf{E} \otimes_{\mathbf{A}} \mathbf{F} \to \mathbf{C} \otimes_{\mathbf{A}} \mathbf{D}\) which, composed with the mapping (22), gives a Z-linear mapping \(g\) of \(\mathbf{E} \otimes_{\mathbf{A}} \mathbf{F}\) into \(\prod_{\lambda, \mu} (\mathbf{E}_{\lambda} \otimes_{\mathbf{A}} F_{\mu})\) such that

\[
g ((x _ {\lambda}) \otimes (y _ {\mu})) = (x _ {\lambda} \otimes y _ {\mu});
\]

moreover, as the families \((x_{\lambda})\) and \((y_{\mu})\) have finite support, so does \((x_{\lambda} \otimes y_{\mu})\) and hence finally g is a canonical homomorphism

\[
g: \left(\bigoplus_ {\lambda \in L} E _ {\lambda}\right) \otimes_ {A} \left(\bigoplus_ {\mu \in M} F _ {\mu}\right)\rightarrow \bigoplus_ {(\lambda , \mu) \in L \times M} (E _ {\lambda} \otimes_ {A} F _ {\mu}),\tag{23}
\]

which is a multimodule homomorphism under the same conditions as (22).

PROPOSITION 7. The canonical mapping (23) is bijective.

To prove this it suffices to define a \(\mathbf{Z}\) -linear mapping \(h\) of the direct sum \(\mathbf{G} = \bigoplus_{(\lambda, \mu) \in L \times M} (\mathbf{E}_{\lambda} \otimes_{\mathbf{A}} \mathbf{F}_{\mu})\) into \(\mathbf{E} \otimes_{\mathbf{A}} \mathbf{F}\) such that \(g \circ h\) and \(h \circ g\) are the identity mappings. But, to define a \(\mathbf{Z}\) -linear mapping of \(\mathbf{G}\) into \(\mathbf{E} \otimes_{\mathbf{A}} \mathbf{F}\) , it suffices (§ 1, no. 6, Proposition 6) to define a \(\mathbf{Z}\) -linear mapping

\[
h _ {\lambda \mu}: \mathrm{E} _ {\lambda} \otimes_ {\mathrm{A}} \mathrm{F} _ {\mu} \rightarrow \mathrm{E} \otimes_ {\mathrm{A}} \mathrm{F}
\]

for every ordered pair \((\lambda, \mu)\) , and we take \(h_{\lambda \mu} = i_{\lambda} \otimes j_{\mu}\) , where \(i_{\lambda}: \mathbf{E}_{\lambda} \to \mathbf{E}\) and \(j_{\mu}: \mathbf{F}_{\mu} \to \mathbf{F}\) are the canonical injections. Then clearly \(h \circ g\) coincides with the identity mapping for the elements of the form \(\left(\sum_{\lambda} x_{\lambda}\right) \otimes \left(\sum_{\mu} y_{\mu}\right)\) which generate the \(\mathbf{Z}\) -module \(\mathbf{E} \otimes_{\mathbf{A}} \mathbf{F}\) ; similarly \(g \circ h\) coincides with the identity mapping for the elements of the form \(\sum_{\lambda, \mu} (x_{\lambda} \otimes y_{\mu})\) which generate the \(\mathbf{Z}\) -module \(\mathbf{G}\) , since for each ordered pair \((\lambda, \mu)\) the products

\[
x _ {\lambda} \otimes y _ {\mu} \quad (x _ {\lambda} \in \mathrm{E} _ {\lambda}, y _ {\mu} \in \mathrm{F} _ {\mu})
\]

generate the Z-module \(E_{\lambda} \otimes_{A} F_{\mu}\) . Whence the proposition.

Let \(u_{\lambda}: \mathbf{E}_{\lambda} \to \mathbf{E}_{\lambda}', v_{\mu}: \mathbf{F}_{\mu} \to \mathbf{F}_{\mu}'\) be A-homomorphisms; clearly the diagram

\[
\begin{array}{c} \left(\bigoplus_ {\lambda} \mathrm{E} _ {\lambda}\right) \otimes_ {\mathrm{A}} \left(\bigoplus_ {\mu} \mathrm{F} _ {\mu}\right) \longrightarrow \bigoplus_ {\lambda , \mu} (\mathrm{E} _ {\lambda} \otimes_ {\mathrm{A}} \mathrm{F} _ {\mu}) \\ (\oplus u _ {\lambda}) \otimes (\oplus v _ {\mu}) \Bigg \downarrow \qquad \qquad \qquad \qquad \qquad \qquad \qquad \qquad \qquad \qquad \qquad \qquad \qquad \qquad \qquad \qquad \qquad \qquad \qquad \oplus (u _ {\lambda} \otimes v _ {\mu}) \\ \left(\bigoplus_ {\lambda} \mathrm{E} _ {\lambda} ^ {\prime}\right) \otimes_ {\mathrm{A}} \left(\bigoplus_ {\mu} \mathrm{F} _ {\mu} ^ {\prime}\right) \longrightarrow \bigoplus_ {\lambda , \mu} (\mathrm{E} _ {\lambda} ^ {\prime} \otimes_ {\mathrm{A}} \mathrm{F} _ {\mu} ^ {\prime}) \end{array}
\]

is commutative.

COROLLARY 1. If the left A-module F admits a basis \((b_{\mu})_{\mu \in \mathbb{M}}\) , every element of \(\mathbf{E} \otimes_{\mathbf{A}} \mathbf{F}\) can be written uniquely in the form \(\sum_{\mu} (x_{\mu} \otimes b_{\mu})\) , where \(x_{\mu} \in \mathbf{E}\) and the family \((x_{\mu})\) has finite support. The \(\mathbf{Z}\) -module \(\mathbf{E} \otimes_{\mathbf{A}} \mathbf{F}\) is isomorphic to \(\mathbf{E}^{(\mathbb{M})}\) considered as a \(\mathbf{Z}\) -module.

The basis \((b_{\mu})\) defines an isomorphism of \(\mathbf{F}\) onto \(\bigoplus_{\mu \in \mathbb{M}}\mathrm{Ab}_{\mu}\) , whence there is an isomorphism \(\mathbf{E} \otimes_{\mathbf{A}}\mathbf{F} \to \bigoplus_{\mu \in \mathbb{M}}(\mathbf{E} \otimes_{\mathbf{A}}\mathbf{Ab}_{\mu})\) by virtue of Proposition 7; as \(\xi \mapsto \xi b_{\mu}\) is an isomorphism of \(\mathbf{A}_s\) onto \(\mathbf{Ab}_{\mu}\) , \(x \mapsto x \otimes b_{\mu}\) is an isomorphism of \(\mathbf{E}\) onto \(\mathbf{E} \otimes_{\mathbf{A}} (\mathbf{Ab}_{\mu})\) by virtue of Proposition 4 of no. 4, whence the corollary. If \(\mathbf{E}\) is a \(((\mathbf{B}_i^{\prime}); \mathbf{A}, (\mathbf{C}_j^{\prime}))\) -multimodule, the canonical isomorphism \(\mathbf{E} \otimes_{\mathbf{A}} \mathbf{F} \to \mathbf{E}^{(\mathbf{M})}\) is a \(((\mathbf{B}_i^{\prime}); (\mathbf{C}_j^{\prime}))\) -multimodule isomorphism.

In particular, if also E admits a basis \((a_{\lambda})_{\lambda \in L}\) , every \(z \in E \otimes_{A} F\) may be written in one and only one way in the form \(\sum_{\lambda, \mu} (a_{\lambda} \xi_{\lambda \mu}) \otimes b_{\mu}\) , where the \(\xi_{\lambda \mu}\) belong to A (and form a family of finite support); the mapping \(z \mapsto (\xi_{\lambda \mu})_{(\lambda, \mu) \in L \times M}\) is an isomorphism of \(E \otimes_{A} F\) onto \(A^{(L \times M)}\) for the Z-module structures (and even the module structures over the centre of A). More particularly:

COROLLARY 2. If E and F are two free modules over a commutative ring C and \((a_{\lambda})\) (resp. \((b_{\mu})\) ) is a basis of the C-module E (resp. F), then \((a_{\lambda} \otimes b_{\mu})\) is a basis of the C-module E \(\otimes_{C} F\) .

By an abuse of language, the basis \((a_{\lambda} \otimes b_{\mu})\) is sometimes called the tensor product of the bases \((a_{\lambda})\) and \((b_{\mu})\) .

Remark (1). Let E be a free right A-module, F a free left A-module, \((a_{\lambda})_{\lambda \in L}\) a basis of E and \((b_{\mu})_{\mu \in M}\) a basis of F. Every element \(z \in E \otimes_{A} F\) may be written uniquely as \(\sum_{\lambda} a_{\lambda} \otimes y_{\lambda}\) , where \(y_{\lambda} \in F\) , and also uniquely as \(\sum_{\mu} x_{\mu} \otimes b_{\mu}\) , where \(x_{\mu} \in E\) . If we write \(y_{\lambda} = \sum_{\mu} \eta_{\lambda \mu} b_{\mu}\) , \(x_{\mu} = \sum_{\lambda} a_{\lambda} \xi_{\lambda \mu}\) , where the \(\xi_{\lambda \mu}\) and \(\eta_{\lambda \mu}\) belong to A, then \(\xi_{\lambda \mu} = \eta_{\lambda \mu}\) for all \((\lambda, \mu)\) , for

\[
\sum_ {\lambda} \left(a _ {\lambda} \otimes \left(\sum_ {\mu} \eta_ {\lambda \mu} b _ {\mu}\right)\right) = \sum_ {\lambda , \mu} \left((a _ {\lambda} \eta_ {\lambda \mu}) \otimes b _ {\mu}\right) = \sum_ {\mu} \left(\left(\sum_ {\lambda} a _ {\lambda} \eta_ {\lambda \mu}\right) \otimes b _ {\mu}\right).
\]

COROLLARY 3. Let \((\mathbf{E}_{\lambda})_{\lambda \in L}\) be a family of right A-modules and F a free (resp. finitely generated free) left A-module. Then the canonical mapping (22)

\[
\left(\prod_ {\lambda \in L} \mathbf {E} _ {\lambda}\right) \otimes_ {\mathbf {A}} \mathbf {F} \rightarrow \prod_ {\lambda \in L} \left(\mathbf {E} _ {\lambda} \otimes_ {\mathbf {A}} \mathbf {F}\right)
\]

is injective (resp. bijective).

If \((b_{\mu})\) is a basis of \(F\) , every element of \(\left(\prod_{\lambda \in L} E_{\lambda}\right) \otimes_{A} F\) can be written uniquely as \(z = \sum_{\mu} ((x_{\lambda}^{(\mu)}) \otimes b_{\mu})\) (Corollary 1); to say that its canonical image is zero means that, for all \(\lambda \in L\) , \(\sum (x_{\lambda}^{(\mu)} \otimes b_{\mu}) = 0\) , hence \(x_{\lambda}^{(\mu)} = 0\) for all \(\lambda \in L\) and all \(\mu\) (Corollary 1) and therefore \(z = 0\) .

Showing that the canonical mapping is bijective when F admits a finite basis is immediately reduced, by virtue of Proposition 7, to the case where

\(F = A_{s}\) ; but then the two sides are canonically identified with \(\prod_{\lambda \in L} E_{\lambda}\) (no. 4, Proposition 4) and after these identifications the canonical mapping (22) becomes the identity.

COROLLARY 4. Let A be a ring with no divisor of zero, E a free right A-module and F a free left A-module. Then the relation \(x \otimes y = 0\) in \(\mathbf{E} \otimes_{\mathbf{A}} \mathbf{F}\) implies \(x = 0\) or \(y = 0\) .

Let \((a_{\lambda})\) be a basis of E, \((b_{\mu})\) a basis of F and let \(x = \sum_{\lambda} a_{\lambda} \xi_{\lambda}, y = \sum_{\mu} \eta_{\mu} b_{\mu}\) ; then \(x \otimes y = \sum_{\lambda, \mu} ((a_{\lambda} \xi_{\lambda} \eta_{\mu}) \otimes b_{\mu})\) and the relation \(x \otimes y = 0\) implies \(\xi_{\lambda} \eta_{\mu} = 0\) for every ordered pairs of indices \((\lambda, \mu)\) (Corollary 1). Hence, if \(x \neq 0\) , that is \(\xi_{\lambda} \neq 0\) for at least one \(\lambda\) , it follows that \(\eta_{\mu} = 0\) for all \(\mu\) , whence \(y = 0\) .

COROLLARY 5. Let E be a right A-module, F a left A-module, M a submodule of E and N a submodule of F. If M is a direct factor of E and N a direct factor of F, the canonical homomorphism \(M \otimes_{A} N \to E \otimes_{A} F\) is injective and the image of \(M \otimes_{A} N\) under this homomorphism is a direct factor of the Z-module \(E \otimes_{A} F\) .

This follows immediately from Proposition 7.

Note that if E is a \(((B_{i}^{\prime}); A, (C_{j}^{\prime}))\) -multimodule and F an \((A, (B_{h}^{\prime\prime}); (C_{k}^{\prime\prime}))\) -multimodule and M and N direct factors in these multimodules, \(M \otimes N\) is a direct factor of the \(((B_{i}^{\prime}), (B_{h}^{\prime\prime}); (C_{j}^{\prime}), (C_{k}^{\prime\prime}))\) -multimodule \(E \otimes F\) .

COROLLARY 6. Let P be a projective left A-module and E, F two right A-modules. For every injective homomorphism \(u: \mathbf{E} \to \mathbf{F}\) , the homomorphism

\[
u \otimes \mathrm{l} _ {\mathbf {P}}: \mathrm{E} \otimes_ {\mathbf {A}} \mathrm{P} \rightarrow \mathrm{F} \otimes_ {\mathbf {A}} \mathrm{P}
\]

is injective.

There exists a left A-module \(\mathbf{Q}\) such that \(\mathbf{L} = \mathbf{P} \oplus \mathbf{Q}\) is free (§2, no. 2, Proposition 4) and \(u \otimes 1_{\mathbf{L}}\) is identified (Proposition 7) with

\[
(u \otimes 1 _ {\mathbf {P}}) \oplus (u \otimes 1 _ {\mathbf {Q}});
\]

it therefore suffices to prove the corollary when P is free (§ 1, no. 6, Corollary 1 to Proposition 7). The same argument reduces the problem to the case where \(P = A_{s}\) , which follows immediately from no. 4, Proposition 4.

COROLLARY 7. Let C be a commutative ring. If E and F are two projective C-modules, E ⊗\_C F is a projective C-module.

This follows immediately from Corollary 5 and the fact that the tensor product of two free C-modules is a free C-module (Corollary 2).

Remark 2. Under the hypotheses of Proposition 7, let \(\mathbf{E}_{\lambda}^{\prime}\) be a submodule of \(\mathbf{E}_{\lambda}, \mathbf{F}_{\mu}^{\prime}\) a submodule of \(\mathbf{F}_{\mu}\) and let \(\mathbf{E}^{\prime} = \bigoplus_{\lambda \in \mathbf{L}} \mathbf{E}_{\lambda}^{\prime}, \mathbf{F}^{\prime} = \bigoplus_{\mu \in \mathbf{M}} \mathbf{F}_{\mu}^{\prime}\) . Let \(\operatorname{Im}(\mathbf{E}^{\prime} \otimes_{\mathbf{A}} \mathbf{F}^{\prime})\)

(resp. \(\operatorname{Im}(\mathbf{E}_{\lambda}^{\prime}\otimes_{\mathbf{A}}\mathbf{F}_{\mu}^{\prime})\) denote the image of \(\mathbf{E}^{\prime}\otimes_{\mathbf{A}}\mathbf{F}^{\prime}\) (resp. \(\mathbf{E}_{\lambda}^{\prime}\otimes_{\mathbf{A}}\mathbf{F}_{\mu}^{\prime}\) ) in \(\mathbf{E}\otimes_{\mathbf{A}}\mathbf{F}\) (resp. \(\mathbf{E}_{\lambda}\otimes_{\mathbf{A}}\mathbf{F}_{\mu}\) ) under the canonical mapping; then the isomorphism (23) identifies the sub- \(\mathbf{Z}\) -modules

\[
\operatorname{Im} \left(\mathrm{E} ^ {\prime} \otimes_ {\mathbf {A}} \mathrm{F} ^ {\prime}\right) \quad \text { and } \quad \bigoplus_ {(\lambda , \mu) \in \mathbf {L} \times \mathbf {M}} \operatorname{Im} \left(\mathrm{E} _ {\lambda} ^ {\prime} \otimes_ {\mathbf {A}} \mathrm{F} _ {\mu} ^ {\prime}\right);
\]

this follows immediately from the commutativity of the diagram

\[
\begin{array}{c} \left(\bigoplus_ {\lambda} \mathrm{E} _ {\lambda} ^ {\prime}\right) \otimes_ {\mathrm{A}} \left(\bigoplus_ {\mu} \mathrm{F} _ {\mu} ^ {\prime}\right) \longrightarrow \left(\bigoplus_ {\lambda} \mathrm{E} _ {\lambda}\right) \otimes_ {\mathrm{A}} \left(\bigoplus_ {\mu} \mathrm{F} _ {\mu}\right) \\ \Biggl \downarrow \\ \bigoplus_ {\lambda , \mu} (\mathrm{E} _ {\lambda} ^ {\prime} \otimes_ {\mathrm{A}} \mathrm{F} _ {\mu} ^ {\prime}) \longrightarrow \bigoplus_ {\lambda , \mu} (\mathrm{E} _ {\lambda} \otimes_ {\mathrm{A}} \mathrm{F} _ {\mu}) \end{array}
\]

where the vertical arrows are the canonical isomorphisms.

\subsection{ASSOCIATIVITY OF THE TENSOR PRODUCT}

PROPOSITION 8. Let A, B be two rings, E a right A-module, F an (A, B)-bimodule and G a left B-module. Then E \(\otimes_{\mathbf{A}} \mathbf{F}\) is a right B-module, F \(\otimes_{\mathbf{B}} \mathbf{G}\) a left A-module and there exists one and only one Z-linear mapping

\[
\phi : (\mathrm{E} \otimes_ {\mathrm{A}} \mathrm{F}) \otimes_ {\mathrm{B}} \mathrm{G} \rightarrow \mathrm{E} \otimes_ {\mathrm{A}} (\mathrm{F} \otimes_ {\mathrm{B}} \mathrm{G})
\]

such that \(\phi((x \otimes y) \otimes z) = x \otimes (y \otimes z)\) for \(x \in \mathbf{E}\) , \(y \in \mathbf{F}\) , \(z \in \mathbf{G}\) ; moreover this \(\mathbf{Z}\) -linear mapping is bijective (``associativity'' of the tensor product).

The right B-module structure on \(\mathbf{E} \otimes_{\mathbf{A}} \mathbf{F}\) and left A-module structure on \(\mathbf{F} \otimes_{\mathbf{B}} \mathbf{G}\) have been defined in no. 4. The uniqueness of \(\phi\) is obvious since the elements \((x \otimes y) \otimes z\) generate the Z-module \((\mathbf{E} \otimes_{\mathbf{A}} \mathbf{F}) \otimes_{\mathbf{B}} \mathbf{G}\) . To show the existence of \(\phi\) , we note that, for all \(z \in \mathbf{G}\) , \(h_z: y \mapsto y \otimes z\) is an A-linear mapping of the left A-module F into the left A-module \(\mathbf{F} \otimes_{\mathbf{B}} \mathbf{G}\) . We write \(g_z = l_{\mathbf{E}} \otimes h_z\) , which is therefore a Z-linear mapping of \(\mathbf{E} \otimes_{\mathbf{A}} \mathbf{F}\) into \(\mathbf{E} \otimes_{\mathbf{A}} (\mathbf{F} \otimes_{\mathbf{B}} \mathbf{G})\) and consider the mapping \((t, z) \mapsto g_z(t)\) from \((\mathbf{E} \otimes_{\mathbf{A}} \mathbf{F} \times \mathbf{G}\) into \(\mathbf{E} \otimes_{\mathbf{A}} (\mathbf{F} \otimes_{\mathbf{B}} \mathbf{G})\) ; as \(h_{z+z'} = h_z + h_{z'}\) for \(z \in \mathbf{G}\) , \(z' \in \mathbf{G}\) , it is immediate that the above mapping is Z-bilinear. Further, we show that for all \(\mu \in \mathbf{B}\) , \(g_{\mu z}(t) = g_z(t\mu)\) ; it is obviously sufficient to do this for \(t = x \otimes y\) where \(x \in \mathbf{E}\) and \(y \in \mathbf{F}\) ; now

\[
g _ {\mu z} (x \otimes y) = x \otimes (y \otimes \mu z)
\]

\[
g _ {z} ((x \otimes y) \mu) = g _ {z} (x \otimes y \mu) = x \otimes (y \mu \otimes z).
\]

Proposition 1 (no. 1) then proves the existence of a Z-linear mapping

\[
\phi : (\mathrm{E} \otimes_ {\mathrm{A}} \mathrm{F}) \otimes_ {\mathrm{B}} \mathrm{G} \rightarrow \mathrm{E} \otimes_ {\mathrm{A}} (\mathrm{F} \otimes_ {\mathrm{B}} \mathrm{G})
\]

such that \(\phi(t \otimes z) = g_z(t)\) , hence \(\phi((x \otimes y) \otimes z) = x \otimes (y \otimes z)\) . Similarly a Z-linear mapping

\[
\psi : \mathbf {E} \otimes_ {\mathbf {A}} (\mathbf {F} \otimes_ {\mathbf {B}} \mathbf {G}) \rightarrow (\mathbf {E} \otimes_ {\mathbf {A}} \mathbf {F}) \otimes_ {\mathbf {B}} \mathbf {G}
\]

is defined such that \(\psi(x \otimes (y \otimes z)) = (x \otimes y) \otimes z\) and clearly \(\psi \circ \phi\) and \(\phi \circ \psi\) are the identity mappings of \((\mathbf{E} \otimes_{\mathbf{A}} \mathbf{F}) \otimes_{\mathbf{B}} \mathbf{G}\) and \(\mathbf{E} \otimes_{\mathbf{A}} (\mathbf{F} \otimes_{\mathbf{B}} \mathbf{G})\) respectively, since they reduce to the identity mappings on generating systems of these \(\mathbf{Z}\) -modules.

It is immediate that, if E is a \(((C_{i}^{\prime}); A, (D_{j}^{\prime}))\) -multimodule, F an \((A, (C_{h}^{\prime\prime}); B, (D_{k}^{\prime\prime}))\) -multimodule and G a \((B, (C_{l}^{\prime\prime}); (D_{m}^{\prime\prime}))\) -multimodule, the canonical isomorphism defined in Proposition 8 is a \(((C_{i}^{\prime}), (C_{h}^{\prime\prime}), (C_{l}^{\prime\prime}); (D_{j}^{\prime}), (D_{k}^{\prime\prime}), (D_{m}^{\prime\prime}))\) -multimodule isomorphism. In particular, if C is a commutative ring and E, F, G three C-modules, there is a canonical C-module isomorphism

\[
(\mathbf {E} \otimes_ {\mathbf {c}} \mathbf {F}) \otimes_ {\mathbf {c}} \mathbf {G} \rightarrow \mathbf {E} \otimes_ {\mathbf {c}} (\mathbf {F} \otimes_ {\mathbf {c}} \mathbf {G}).
\]

We shall see below that, under certain conditions, the definition of tensor product can be generalized to a family of multimodules, which will in particular give us under the hypotheses of Proposition 8 a Z-module \(E \otimes_{A} F \otimes_{B} G\) , which is canonically isomorphic to each of the Z-modules \((E \otimes_{A} F) \otimes_{B} G\) and \(E \otimes_{A} (F \otimes_{B} G)\) and with which the latter will be identified.

\subsection{TENSOR PRODUCT OF FAMILIES OF MULTIMODULES}

Let \((\mathbf{G}_{\lambda})_{\lambda \in L}\) be a family of \(\mathbf{Z}\) -modules; a mapping \(u\) of the set \(\mathbf{G} = \prod_{\lambda \in L} \mathbf{G}_{\lambda}\) into a \(\mathbf{Z}\) -module is said to be multiadditive (or \(\mathbf{Z}\) -multilinear) if \((x_{\lambda}) \mapsto u((x_{\lambda}))\) is additive with respect to each of the variables \(x_{\lambda}\) ; to be precise, this means that, for all \(\mu \in L\) and every element \((a_{\lambda}) \in \prod_{\lambda \neq \mu} \mathbf{G}_{\lambda}\) , canonically identifying \(\mathbf{G}\) with \(\mathbf{G}_{\mu} \times \prod_{\lambda \neq \mu} \mathbf{G}_{\lambda}\) ,

\[
u (x _ {\mu} + y _ {\mu}, (a _ {\lambda})) = u (x _ {\mu}, (a _ {\lambda})) + u (y _ {\mu}, (a _ {\lambda})) \quad \text { for } x _ {\mu}, y _ {\mu} \text { in } G _ {\mu}.\tag{24}
\]

This implies in particular that \(u((x_{\lambda})) = 0\) if one of the \(x_{\lambda}\) is zero.

We consider also the universal mapping problem where \(\Sigma\) is the species of Z-module structure and the \(\alpha\) -mappings are the multiadditive mappings of G into a Z-module. A solution is still obtained by considering the Z-module \(C = Z^{(G)}\) of formal linear combinations of elements of G with coefficients in Z and the sub-Z-module D of C generated by the elements of the form

\[
\left(x _ {\mu} + y _ {\mu}, \left(z _ {\lambda}\right) _ {\lambda \neq \mu}\right) - \left(x _ {\mu}, \left(z _ {\lambda}\right) _ {\lambda \neq \mu}\right) - \left(y _ {\mu}, \left(z _ {\lambda}\right) _ {\lambda \neq \mu}\right)
\]

where \(\mu \in \mathbf{L}\) , \(x_{\mu} \in \mathrm{G}_{\mu}\) , \(y_{\mu} \in \mathrm{G}_{\mu}\) and the \(z_{\lambda} \in \mathrm{G}_{\lambda}\) ( \(\lambda \neq \mu\) ) are arbitrary. The quotient \(\mathbf{Z}\) -module C/D is called the tensor product (over \(\mathbf{Z}\) ) of the family \((\mathrm{G}_{\lambda})_{\lambda \in L}\) of \(\mathbf{Z}\) -modules and is denoted by \(\bigotimes_{\lambda \in L} \mathrm{G}_{\lambda}\) ; for every element \((x_{\lambda})_{\lambda \in L}\) of G which is an element of the canonical basis of C, \(\bigotimes_{\lambda \in L} x_{\lambda}\) denotes the canonical image of this element in C/D. It follows from the above definitions that the mapping \(\phi: (x_{\lambda}) \mapsto \bigotimes_{\lambda \in L} x_{\lambda}\) of G into \(\bigotimes_{\lambda \in L} G_{\lambda}\) is Z-multilinear and that, for every Z-multilinear mapping \(f\) of G into a Z-module H, there exists one and only one Z-linear mapping \(g: \bigotimes_{\lambda \in L} G_{\lambda} \to H\) such that \(f = g \circ \phi\) ; the ordered pair \(\left( \bigotimes_{\lambda \in L} G_{\lambda}, \phi \right)\) thus resolves the universal mapping problem in question.

Let \((\mathbf{G}_{\lambda}^{\prime})_{\lambda \in L}\) be another family of \(\mathbf{Z}\) -modules and, for all \(\lambda \in L\) , let \(v_{\lambda}: \mathbf{G}_{\lambda} \to \mathbf{G}_{\lambda}^{\prime}\) be a \(\mathbf{Z}\) -linear mapping (in other words a homomorphism of commutative groups). Then the mapping

\[
(x _ {\lambda}) \mapsto \bigotimes_ {\lambda \in L} v _ {\lambda} (x _ {\lambda})
\]

of \(\mathbf{G}\) into \(\bigotimes_{\lambda \in L} \mathbf{G}_{\lambda}'\) is \(\mathbf{Z}\) -multilinear and hence defines canonically a \(\mathbf{Z}\) -linear mapping of \(\bigotimes_{\lambda \in L} \mathbf{G}_{\lambda}\) into \(\bigotimes_{\lambda \in L} \mathbf{G}_{\lambda}'\) denoted by \(\bigotimes_{\lambda \in L} v_{\lambda}\) and such that

\[
\left(\bigotimes_ {\lambda \in L} v _ {\lambda}\right) \left(\bigotimes_ {\lambda \in L} x _ {\lambda}\right) = \bigotimes_ {\lambda \in L} v _ {\lambda} (x _ {\lambda}).\tag{25}
\]

In particular, we consider, for some \(\mu \in \mathbf{L}\) , an endomorphism \(\theta\) of \(\mathbf{G}_{\mu}\) ; we denote by \(\tilde{\theta}\) the endomorphism of \(\bigotimes_{\lambda \in \mathbf{L}} \mathbf{G}_{\lambda}\) equal to \(\bigotimes_{\lambda \in \mathbf{L}} v_{\lambda}\) where \(v_{\mu} = \theta\) and \(v_{\lambda} = 1_{\mathbf{G}_{\lambda}}\) for \(\lambda \neq \mu\) .

Then we suppose given a set \(\Omega\) , a mapping

\[
c: \omega \mapsto (\rho (\omega), \sigma (\omega))
\]

of \(\Omega\) into \(\mathbf{L} \times \mathbf{L}\) and, for all \(\omega \in \Omega\) , an endomorphism \(p_{\omega}\) of \(\mathrm{G}_{\rho(\omega)}\) and an endomorphism \(q_{\omega}\) of \(\mathrm{G}_{\sigma(\omega)}\) ; there correspond to them two endomorphisms \(\tilde{p}_{\omega}\) and \(\tilde{q}_{\omega}\) of \(\mathrm{P} = \bigotimes_{\lambda \in \mathbf{L}} \mathrm{G}_{\lambda}\) . Let \(\mathbf{R}\) be the sub- \(\mathbf{Z}\) -module of \(\mathbf{P}\) generated by the union of the images of the endomorphisms \(\tilde{p}_{\omega} - \tilde{q}_{\omega}\) when \(\omega\) runs through \(\Omega\) . The quotient \(\mathbf{Z}\) -module \(\mathrm{P/R}\) is called the tensor product of the family \((\mathrm{G}_{\lambda})_{\lambda \in \mathbf{L}}\) relative to \(c, p, q\) and is denoted by \(\bigotimes_{(c, p, q)} \mathrm{G}_{\lambda}\) ; composing the canonical homomorphism \(\mathrm{P} \to \mathrm{P/R}\) with the mapping \(\phi: \mathrm{G} \to \bigotimes_{\lambda \in \mathbf{L}} \mathrm{G}_{\lambda}\) defined above, we obtain a \(\mathbf{Z}\) -multilinear mapping \(\phi_{(c, p, q)}: \mathrm{G} \to \bigotimes_{(c, p, q)} \mathrm{G}_{\lambda}\) and write \(\phi_{(c, p, q)}((x_{\lambda})) = \bigotimes_{(c, p, q)} x_{\lambda}\) or simply \(\bigotimes_{(c)} x_{\lambda}\) . The ordered pair consisting of \(\bigotimes_{(c, p, q)} x_{\lambda}\) and \(\phi_{(c, p, q)}\) resolves the following universal mapping problem: let \(\tilde{p}_{\omega}\) (resp. \(\tilde{q}_{\omega}\) ) denote the mapping

\[
\left(x _ {\rho (\omega)}, \left(x _ {\lambda}\right) _ {\lambda \neq \rho (\omega)}\right) \mapsto \left(p _ {\omega} \left(x _ {\rho (\omega)}\right), \left(x _ {\lambda}\right) _ {\lambda \neq \rho (\omega)}\right)
\]

\[
\left(\text { resp. } \left(x _ {\sigma (\omega)}, \left(x _ {\lambda}\right) _ {\lambda \neq \sigma (\omega)}\right) \mapsto \left(q _ {\omega} (x _ {\sigma (\omega)}), \left(x _ {\lambda}\right) _ {\lambda \neq \sigma (\omega)}\right)\right)
\]

of G into itself. Then \(\Sigma\) is taken to be the species of Z-module structure and the \(\alpha\) -mappings the \(\mathbf{Z}\) -multilinear mappings \(u\) of \(\mathbf{G}\) into a \(\mathbf{Z}\) -module which further satisfy the conditions

\[
u \circ \tilde {p} _ {\omega} = u \circ \tilde {q} _ {\omega}\tag{26}
\]

for all \(\omega \in \Omega\) . The proof is obvious from the above definitions.

This construction recovers in particular that of \(\mathbf{E} \otimes_{\mathbf{A}} \mathbf{F}\) described in no. 1: in this case we take \(\mathbf{L} = \{1, 2\}\) , \(\mathbf{G}_1 = \mathbf{E}\) , \(\mathbf{G}_2 = \mathbf{F}\) , \(\Omega = \mathbf{A}\) ; further, for all \(\omega \in \mathbf{A}\) , we must have \(\rho(\omega) = 1\) , \(\sigma(\omega) = 2\) , \(p_\omega\) is the endomorphism \(x \mapsto x\omega\) of the \(\mathbf{Z}\) -module \(\mathbf{E}\) and \(q_\omega\) the endomorphism \(y \mapsto \omega y\) of the \(\mathbf{Z}\) -module \(\mathbf{F}\) .

Let \((\mathbf{G}_{\lambda}^{\prime})_{\lambda \in L}\) be a second family of Z-modules; keeping the mapping c the same, suppose given for all \(\omega \in \Omega\) , an endomorphism \(p_{\omega}^{\prime}\) of \(\mathbf{G}_{\rho(\omega)}^{\prime}\) and an endomorphism \(q_{\omega}^{\prime}\) of \(\mathbf{G}_{\sigma(\omega)}^{\prime}\) . For all \(\lambda \in L\) , then let \(v_{\lambda}: G_{\lambda} \to G_{\lambda}^{\prime}\) be a Z-linear mapping such that, for all \(\omega \in \Omega\) ,

\[
v _ {\rho (\omega)} \circ p _ {\omega} = p _ {\omega} ^ {\prime} \circ v _ {\rho (\omega)} \quad \text { and } \quad v _ {\sigma (\omega)} \circ q _ {\omega} = q _ {\omega} ^ {\prime} \circ v _ {\sigma (\omega)}\tag{27}
\]

(in other words, for all \(\lambda \in \mathbf{L}\) , \(v_{\lambda}\) is a morphism for the laws of action on \(\mathbf{G}_{\lambda}\) (resp. \(\mathrm{G}_{\lambda}^{\prime}\) ) defined on the \(p_{\xi}\) and \(q_{\eta}\) (resp. \(p_{\xi}^{\prime}\) and \(q_{\eta}^{\prime}\) ) with \(\xi\) and \(\eta\) such that \(\rho(\xi) = \lambda\) and \(\sigma(\eta) = \lambda\) ). Then the mapping

\[
u: (x _ {\lambda}) \mapsto \bigotimes_ {(c, p', q')} v _ {\lambda} (x _ {\lambda})
\]

of G into \(\bigotimes_{(c, p', q')} \mathrm{G}_{\lambda}'\) is Z-multilinear and satisfies conditions (26) and hence defines a Z-linear mapping of \(\bigotimes_{(c, p, q)} \mathrm{G}_{\lambda}\) into \(\bigotimes_{(c, p', q')} \mathrm{G}_{\lambda}'\) , which we shall denote simply by \(\bigotimes_{(c)} v_{\lambda}\) if no confusion can arise.

We shall now give an ``associativity'' property for the general tensor products thus defined. Let \((\mathbf{L}_{i})_{1\leq i\leq n}\) be a finite partition of L; for every index i, let \(\Omega_{i}\) denote the subset of \(\Omega\) consisting of the elements such that \(\rho(\omega)\in\mathbf{L}_{i}\) and \(\sigma(\omega)\in\mathbf{L}_{i}\) ; clearly the \(\Omega_{i}\) are pairwise disjoint; we set \(\Omega^{\prime}=\Omega-\left(\bigcup_{i}\Omega_{i}\right)\) . For each index i, \(c^{(i)}\) will denote the mapping \(\omega\mapsto(\rho(\omega),\sigma(\omega))\) of \(\Omega_{i}\) into \(L_{i}\times L_{i}\) ; for \(\omega\in\Omega_{i}\) , we write \(p_{\omega}^{(i)}\) and \(q_{\omega}^{(i)}\) instead of \(p_{\omega}\) and \(q_{\omega}\) . Then for each i there is a ``partial'' tensor product

\[
\mathrm{F} _ {i} = \bigotimes_ {(c ^ {(i)}, p ^ {(i)}, q ^ {(i)})} \mathrm{G} _ {\lambda}.
\]

We shall further make the following ``permutability'' hypothesis:

\begin{enumerate}
\def\labelenumi{(\Alph{enumi})}
\setcounter{enumi}{15}
\tightlist
\item
  If \(\omega \in \Omega'\) , \(p_{\omega}\) (resp. \(q_{\omega}\) ) permutes with each of the endomorphisms \(p_{\xi}\) and \(q_{\eta}\) of \(\mathrm{G}_{\rho(\omega)}\) (resp. \(\mathrm{G}_{\sigma(\omega)}\) ) such that \(\xi \notin \Omega'\) , \(\eta \notin \Omega'\) and \(\rho(\omega) = \rho(\xi) = \sigma(\eta)\) (resp. \(\sigma(\omega) = \rho(\xi) = \sigma(\eta)\) ).
\end{enumerate}

For each \(\omega \in \Omega'\) , let \(i\) be the index such that \(\rho(\omega) \in \mathbf{L}_i\) ; then consider the family \((v_{\lambda})_{\lambda \in \mathbf{L}_i}\) where \(v_{\rho(\omega)} = p_{\omega}\) and \(v_{\lambda} = 1_{\mathbf{G}_{\lambda}}\) for \(\lambda \neq \rho(\omega)\) ; hypothesis (P) implies that the family \((v_{\lambda})\) satisfies conditions (27) (where \(p'\) and \(p\) must be replaced by \(p^{(i)}\) , \(q'\) and \(q\) by \(q^{(i)}\) , \(\omega\) by an element \(\xi\) running through \(\Omega_i\) ); thus an endomorphism \(\bigotimes_{(c^{(i)})} v_{\lambda} = r_{\omega}\) is derived of the \(\mathbf{Z}\) -module \(\mathbf{F}_i\) . Similarly, an endomorphism \(s_{\omega}\) is defined of the \(\mathbf{Z}\) -module \(\mathbf{F}_j\) starting with \(q_{\omega}, j\) being the index such that \(\sigma(\omega) \in \mathbf{L}_j\) ; finally let \(d(\omega) = (i,j)\) . Then we can define the tensor product \(\bigotimes_{(d,r,s)} \mathbf{F}_i\) and the corresponding canonical mapping

\[
\phi_ {(d, r, s)}: \prod_ {i = 1} ^ {n} \mathrm{F} _ {i} \rightarrow \bigotimes_ {(d, r, s)} \mathrm{F} _ {i}.
\]

On the other hand, for each i, the canonical mapping

\[
\psi_ {i} = \phi_ {(c ^ {(i)}, p ^ {(i)}, q ^ {(i)})}: \prod_ {\lambda \in L _ {i}} G _ {\lambda} \rightarrow F _ {i};
\]

using the associativity of the product of sets, a Z-multilinear mapping \(\psi = \phi_{(d,r,s)}\circ (\psi_i)\) of G into \(\bigotimes_{(d,r,s)}\mathrm{F}_i\) is derived. We show that the ordered pair \(\left(\bigotimes_{(d,r,s)}\mathrm{F}_i,\psi\right)\) is a solution of the same universal problem as \(\left(\bigotimes_{(c,p,q)}\mathrm{G}_{\lambda},\phi_{(c,p,q)}\right)\) , whence will follow the existence of a unique Z-module isomorphism

\[
\theta \colon \bigotimes_ {(c, p, q)} \mathrm{G} _ {\lambda} \rightarrow \bigotimes_ {(d, r, s)} \mathrm{F} _ {i}
\]

such that \(\psi = \theta \circ \phi_{(c,p,q)}\) (Set Theory, IV, § 3, no. 1).

By induction on \(n\) , the proof is reduced to the case \(n = 2\) ; for simplicity we write \(\mathbf{F}_1 \otimes_{(d)} \mathbf{F}_2\) and \(y_1 \otimes_{(d)} y_2\) instead of \(\bigotimes_{(d, r, s)} \mathbf{F}_i\) and \(\bigotimes_{(d, r, s)} y_i\) . Consider the mapping from \(G\) to \(\mathbf{F}_1 \otimes_{(d)} \mathbf{F}_2\)

\[
h \colon (x _ {\lambda}) \to \left(\bigotimes_ {(c ^ {(1)})} x _ {\lambda}\right) \underset {(d)} {\otimes} \left(\bigotimes_ {(c ^ {(2)})} x _ {\lambda}\right).
\]

It is obviously Z-multilinear; we show that it satisfies conditions (26) for all \(\omega\in\Omega\) . This is obvious if \(\omega\in\Omega_{1}\) or \(\omega\in\Omega_{2}\) ; otherwise, supposing, in order to fix the ideas, that \(\rho(\omega)\in\mathbf{L}_{1}\) and \(\sigma(\omega)\in\mathbf{L}_{2}\) , the values of \(h\circ\tilde{p}_{\omega}\) and \(h\circ\tilde{q}_{\omega}\) at \((x_{\lambda})\) are respectively

\[
\left(r _ {\omega} \left(\bigotimes_ {(c ^ {(1)})} x _ {\lambda}\right)\right) \underset {(d)} {\otimes} \left(\bigotimes_ {(c ^ {(2)})} x _ {\lambda}\right) \quad \text { and } \quad \left(\bigotimes_ {(c ^ {(1)})} x _ {\lambda}\right) \underset {(d)} {\otimes} \left(s _ {\omega} \left(\bigotimes_ {(c ^ {(2)})} x _ {\lambda}\right)\right)
\]

which are also equal by definition of \(\mathbf{F}_1\otimes_{(d)}\mathbf{F}_2\)

This being so, let \(u\) be a \(\mathbf{Z}\) -multilinear mapping of \(G\) into a \(\mathbf{Z}\) -module \(H\) , satisfying conditions (26); we shall define a \(\mathbf{Z}\) -linear mapping \(v: F_1 \otimes F_2 \to H\) such that \(u = v \circ h\) and that will prove our assertion (repeating the argument of no. 1, Corollary 1 to Proposition 1). For all \(z_2 = (x_\lambda)_{\lambda \in L_2}\) we consider the ``partial'' linear mapping of \(\prod_{\lambda \in L_1} G_\lambda\) into \(H\)

\[
u \left(\cdot , z _ {2}\right): \left(x _ {\lambda}\right) _ {\lambda \in L _ {1}} \mapsto u \left(\left(x _ {\lambda}\right) _ {\lambda \in L _ {1}}, z _ {2}\right) = u \left(\left(x _ {\lambda}\right) _ {\lambda \in L}\right).\tag{28}
\]

Clearly it is Z-multilinear and satisfies conditions (26) for \(\omega\in\Omega_{1}\) ; by definition there thus exists a Z-linear mapping \(y_{1}\mapsto w_{1}(y_{1},x_{2})\) of \(F_{1}\) into H such that

\[
w _ {1} \left(\bigotimes_ {(c ^ {(1)})} x _ {\lambda}, z _ {2}\right) = u \left(\left(x _ {\lambda}\right) _ {\lambda \in L _ {1}}, z _ {2}\right)\tag{29}
\]

We next consider the mapping

\[
u _ {2}: (x _ {\lambda}) _ {\lambda \in L _ {2}} \mapsto w _ {1} (\cdot , (x _ {\lambda}) _ {\lambda \in L _ {2}})
\]

of \(\prod_{\lambda \in L_2} G_\lambda\) into \(\mathrm{Hom}_{\mathbf{Z}}(\mathrm{F}_1, \mathrm{H})\) ; it is obviously \(\mathbf{Z}\) -multilinear and satisfies conditions (26) for \(\omega \in \Omega_2\) , by virtue of the hypothesis on \(u\) and relations (28) and (29) and taking account of the fact that the elements of the form \(\bigotimes_{(c^{(1)})} x_\lambda\) generate the \(\mathbf{Z}\) -module \(\mathrm{F}_1\) . Hence there is a \(\mathbf{Z}\) -linear mapping

\[
w _ {2}: \mathrm{F} _ {2} \rightarrow \operatorname{Hom} _ {\mathbf {Z}} (\mathrm{F} _ {1}, \mathrm{H})
\]

such that

\[
w _ {2} \left(\bigotimes_ {(c ^ {(2)})} x _ {\lambda}\right) = u _ {2} \left(\left(x _ {\lambda}\right) _ {\lambda \in L _ {2}}\right)
\]

or also

\[
\left(w _ {2} \left(\bigotimes_ {(c ^ {(2)})} x _ {\lambda}\right)\right) \left(\bigotimes_ {(c ^ {(1)})} x _ {\lambda}\right) = u \left(\left(x _ {\lambda}\right) _ {\lambda \in L}\right).\tag{30}
\]

We now consider, for \(y_{1} \in F_{1}, y_{2} \in F_{2}\) , the element of H

\[
w (y _ {1}, y _ {2}) = (w _ {2} (y _ {2})) (y _ {1}).\tag{31}
\]

Clearly \(w\) is a \(\mathbf{Z}\) -bilinear mapping of \(\mathbf{F}_1 \times \mathbf{F}_2\) into \(H\) . We show further that, for all \(\omega \in \Omega'\) , (assuming, to fix the ideas, that \(\rho(\omega) \in L_1\) and \(\sigma(\omega) \in L_2\) )

\[
w (r _ {\omega} (y _ {1}), y _ {2}) = w (y _ {1}, s _ {\omega} (y _ {2})).\tag{32}
\]

It suffices to verify this relation when \(y_{1}\) (resp. \(y_{2}\) ) is of the form \(\bigotimes_{(c^{(1)})} x_{\lambda}\) (resp. \(\bigotimes_{(c^{(2)})} x_{\lambda}\) ), since these elements generate the Z-module \(F_{1}\) (resp. \(F_{2}\) ). But by definition, \(r_{\omega}\left(\bigotimes_{(c^{(1)})} x_{\lambda}\right) = \bigotimes_{(c^{(1)})} x_{\lambda}'\) , where \(x_{\rho(\omega)}' = p_{\omega}(x_{\rho(\omega)})\) and \(x_{\lambda}' = x_{\lambda}\) for \(\lambda \neq \rho (\omega)\) in \(\mathbf{L}_1\) ; similarly \(s_{\omega}\left(\bigotimes_{(c^{(2)})}x_{\lambda}\right) = \bigotimes_{(c^{(2)})}x_{\lambda}^{\prime \prime}\) , where \(x_{\sigma (\omega)}^{\prime \prime} = q_{\omega}(x_{\sigma (\omega)})\) and \(x_{\lambda}^{\prime \prime} = x_{\lambda}\) for \(\lambda \neq \sigma (\omega)\) in \(\mathbf{L}_2\) ; using (30) and (31), relation (32) then follows from (26). Hence there exists a Z-linear mapping \(v\) of \(\mathrm{F}_1\otimes \mathrm{F}_2\) into H such that \(v(y_{1}\otimes y_{2}) = w(y_{1},y_{2})\) and it then follows from (30) and (31) that \(v\circ h = u\) .

The most important special case of the general tensor product defined above is the following: we start with a family \((\mathbf{A}_i)_{1\leqslant i\leqslant n - 1}\) of rings and a family \((\mathbf{E}_i)_{1\leqslant i\leqslant n}\) , where \(\mathbf{E}_1\) is a right \(\mathbf{A}_1\) -module, \(\mathbf{E}_n\) is a left \(\mathbf{A}_{n - 1}\) module and for \(2\leqslant i\leqslant n - 1\) , \(\mathbf{E}_i\) is an \((\mathbf{A}_{i - 1},\mathbf{A}_i)\) -bimodule. Then the above definition is applied as follows: L is the set \([1,n]\) , \(\mathrm{G}_i = \mathrm{E}_i\) , \(\Omega\) is the set the sum of the \(\mathbf{A}_i\) ( \(1\leqslant i\leqslant n - 1\) ). For \(\omega \in \mathbf{A}_i\) ( \(1\leqslant i\leqslant n - 1\) ), take \(\rho (\omega) = i\) , \(\sigma (\omega) = i + 1\) , \(p_{\omega}\) is the endomorphism \(x\mapsto x\omega\) of the Z-module \(\mathbf{E}_i\) and \(q_{\omega}\) the endomorphism \(y\mapsto \omega y\) of the Z-module \(\mathbf{E}_{i + 1}\) ; the corresponding tensor product is denoted by

\[
\mathbf {E} _ {1} \otimes_ {\mathbf {A} _ {1}} \mathbf {E} _ {2} \otimes_ {\mathbf {A} _ {2}} \mathbf {E} _ {3} \otimes \dots \otimes_ {\mathbf {A} _ {n - 2}} \mathbf {E} _ {n - 1} \otimes_ {\mathbf {A} _ {n - 1}} \mathbf {E} _ {n}\tag{33}
\]

(a notation where the \(A_{i}\) may occasionally be suppressed) and the elements \(\bigotimes_{(c,p,q)} x_{i}\) of this tensor product, for a family \((x_{i})\) such that \(x_{i} \in E_{i}\) for \(1 \leqslant i \leqslant n\) , may be written \(x_{1} \otimes x_{2} \otimes \cdots \otimes x_{n}\) if no confusion can arise; an analogous notation is used for a Z-linear mapping \(\bigotimes_{(c)} v_{i}\) . Hypothesis (P) holds for every partition of \([1,n]\) , since the \(E_{i}\) are bimodules for \(2 \leqslant i \leqslant n - 1\) . When n = 3, we have thus defined the Z-module \(E \otimes_{A} F \otimes_{B} G\) alluded to in no. 8, and recovered Proposition 8 (no. 8).

When each of the \(E_{i}\) is a multimodule (with, for \(2 \leqslant i \leqslant n - 1\) , \(A_{i-1}\) one of the rings operating on the left and \(A_{i}\) one of the rings operating on the right and analogous conditions for i = 1 and i = n), as in no. 4, a multimodule structure is defined on \(E_{1} \otimes_{A_1} E_{2} \otimes \cdots \otimes_{A_{n-1}} E_{n}\) with respect to all the rings except the \(A_{i}\) which operate on the \(E_{i}\) ( \(1 \leqslant i \leqslant n\) ).

In particular, let C be a commutative ring, \((\mathbf{E}_{i})_{1\leqslant i\leqslant n}\) a family of C-modules. By giving \(E_{1}\) and \(E_{n}\) two C-module structures identical with the given structure and \(E_{i}\) for \(2\leqslant i\leqslant n-1\) three C-module structures identical with the given structure, we define on the tensor product

\[
\mathbf {E} _ {1} \otimes_ {\mathbf {c}} \mathbf {E} _ {2} \otimes_ {\mathbf {c}} \mathbf {E} _ {3} \otimes \dots \otimes_ {\mathbf {c}} \mathbf {E} _ {n - 1} \otimes_ {\mathbf {c}} \mathbf {E} _ {n}\tag{34}
\]

\(n\) C-modules structures which are compatible with one another and which are in fact identical, since for all \(\gamma \in \mathbf{C}\) and \((x_{i})\in \prod_{i = 1}^{n}\mathbf{E}_{i}\) , by definition

\[
\begin{array}{r l} \left(\gamma x _ {1}\right) \otimes x _ {2} \otimes \dots \otimes x _ {n} \\ & = x _ {1} \otimes \left(\gamma x _ {2}\right) \otimes \dots \otimes x _ {n} = \dots = x _ {1} \otimes x _ {2} \otimes \dots \otimes \left(\gamma x _ {n}\right). \end{array}
\]

When we speak of the tensor product (34) as a C-module, we shall always mean with this structure, unless otherwise mentioned, and the tensor product (34) is also denoted by \(\bigotimes_{1\leqslant i\leqslant n}\mathbf{E}_i\) if no confusion can arise. For every C-module G, the Z-multilinear mappings of \(\prod_{i=1}^{n}\mathbf{E}_i\) into G which, for every index \(i\) , satisfy the relation

\[
f (x _ {1}, \dots , x _ {i - 1}, \gamma x _ {i}, x _ {i + 1}, \dots , x _ {n}) = \gamma f (x _ {1}, \dots , x _ {n})\tag{35}
\]

for \(\gamma \in \mathbf{C}\) and \((x_{i})\in \prod_{i}\mathbf{E}_{i}\) are then called C-multilinear and form a C-module denoted by \(\mathcal{L}_n(\mathbf{E}_1,\dots ,\mathbf{E}_n;\mathbf{G})\) ; the universal property of the tensor product (34) then allows us to define a canonical C-module isomorphism

\[
\mathcal {L} _ {n} \left(\mathrm{E} _ {1}, \dots , \mathrm{E} _ {n}; \mathrm{G}\right)\rightarrow \operatorname{Hom} _ {\mathrm{C}} \left(\mathrm{E} _ {1} \otimes_ {\mathrm{C}} \mathrm{E} _ {2} \otimes \dots \otimes_ {\mathrm{C}} \mathrm{E} _ {n}, \mathrm{G}\right)\tag{36}
\]

which associates with every C-multilinear mapping f the C-linear mapping g such that

\[
f (x _ {1}, \dots , x _ {n}) = g (x _ {1} \otimes x _ {2} \otimes \dots \otimes x _ {n}).
\]

A C-multilinear mapping of \(E_{1} \times \cdots \times E_{n}\) into C is also called an n-linear form.

Let \((\mathbf{F}_i)_{1\leqslant i\leqslant n}\) be another family of C-modules; for every system of \(n\) C-linear mappings \(u_i: \mathbf{E}_i \to \mathbf{F}_i\) , \(u_1 \otimes u_2 \otimes \cdots \otimes u_n\) (also denoted by \(\bigotimes_{1\leqslant i\leqslant n} u_i\) ) is a C-linear mapping of

\[
\mathbf {E} _ {1} \otimes_ {\mathbb {C}} \mathbf {E} _ {2} \otimes \dots \otimes_ {\mathbb {C}} \mathbf {E} _ {n} \quad \text { into } \mathbf {F} _ {1} \otimes_ {\mathbb {C}} \mathbf {F} _ {2} \otimes \dots \otimes_ {\mathbb {C}} \mathbf {F} _ {n}.
\]

Further, \((u_{1},\ldots,u_{n})\mapsto u_{1}\otimes u_{2}\otimes\cdots\otimes u_{n}\) is a C-multilinear mapping of \(\prod_{i}\operatorname{Hom}_{\mathbb{C}}(\mathbf{E}_{i},\mathbf{F}_{i})\) into

\[
\operatorname{Hom} _ {\mathbf {c}} (\mathbf {E} _ {1} \otimes_ {\mathbf {c}} \mathbf {E} _ {2} \otimes \dots \otimes_ {\mathbf {c}} \mathbf {E} _ {n}, \mathbf {F} _ {1} \otimes_ {\mathbf {c}} \mathbf {F} _ {2} \otimes \dots \otimes_ {\mathbf {c}} \mathbf {F} _ {n}).
\]

Hence there corresponds canonically to the latter mapping a C-linear mapping called canonical

\[
\begin{array}{r l} & \operatorname{Hom} _ {\mathbf {c}} (\mathrm{E} _ {1}, \mathrm{F} _ {1}) \otimes_ {\mathbf {c}} \operatorname{Hom} _ {\mathbf {c}} (\mathrm{E} _ {2}, \mathrm{F} _ {2}) \otimes \dots \otimes_ {\mathbf {c}} \operatorname{Hom} _ {\mathbf {c}} (\mathrm{E} _ {n}, \mathrm{F} _ {n}) \\ & \qquad \to \operatorname{Hom} _ {\mathbf {c}} (\mathrm{E} _ {1} \otimes_ {\mathbf {c}} \mathrm{E} _ {2} \otimes \dots \otimes_ {\mathbf {c}} \mathrm{E} _ {n}, \mathrm{F} _ {1} \otimes_ {\mathbf {c}} \mathrm{F} _ {2} \otimes \dots \otimes_ {\mathbf {c}} \mathrm{F} _ {n}) \end{array}\tag{37}
\]

generalizing that defined in no. 5 for \(n = 2\) .

The general associativity property seen earlier can be specialized here as follows. Given a partition \((\mathbf{J}_{k})_{1\leqslant k\leqslant m}\) of the interval \([1,n]\) of N, let, for each k, \(F_{k}\) be the tensor product \(E_{i_{1}}\otimes_{C}E_{i_{2}}\otimes\cdots\otimes_{C}E_{i_{r}}\) , where \((i_{1},\ldots,i_{r})\) is the strictly increasing sequence of elements of \(J_{k}\) . Then there is a canonical C-module isomorphism (called ``associativity isomorphism'')

\[
\mathrm{F} _ {1} \otimes_ {\mathbf {c}} \mathrm{F} _ {2} \otimes \dots \otimes_ {\mathbf {c}} \mathrm{F} _ {m} \rightarrow \mathrm{E} _ {1} \otimes_ {\mathbf {c}} \mathrm{E} _ {2} \otimes \dots \otimes_ {\mathbf {c}} \mathrm{E} _ {n}
\]

which, in the above notation, maps the tensor product

\[
y _ {1} \otimes y _ {2} \otimes \dots \otimes y _ {m}, \quad \text { where } y _ {k} = x _ {i _ {1}} \otimes x _ {i _ {2}} \otimes \dots \otimes x _ {i _ {r}},
\]

to the tensor product \(x_{1} \otimes x_{2} \otimes \cdots \otimes x_{n}\) (where \(x_{i} \in E_{i}\) for all i).

In particular, if \(\pi\) is a permutation of \([1, n]\) , writing \(\mathbf{J}_k = \{\pi(k)\}\) for \(1 \leqslant k \leqslant n\) , we obtain a canonical (``commutativity'') isomorphism

\[
\mathbf {E} _ {\pi (1)} \otimes_ {\mathbf {C}} \mathbf {E} _ {\pi (2)} \otimes \dots \otimes_ {\mathbf {C}} \mathbf {E} _ {\pi (n)} \rightarrow \mathbf {E} _ {1} \otimes_ {\mathbf {C}} \mathbf {E} _ {2} \otimes \dots \otimes_ {\mathbf {C}} \mathbf {E} _ {n}
\]

which maps \(x_{\pi(1)} \otimes x_{\pi(2)} \otimes \cdots \otimes x_{\pi(n)}\) to \(x_{1} \otimes x_{2} \otimes \cdots \otimes x_{n}\) . We shall often identify the various tensor products which correspond to one another under these canonical isomorphisms.

For \(1 \leqslant i \leqslant n\) , suppose that \(\mathbf{E}_i\) admits a basis \((b_{\lambda_1}^{(i)})_{\lambda_i \in \mathbf{L}_i}\) ; by induction on \(n\) , it follows from no. 7, Corollary 2 to Proposition 7 that the family \((b_{\lambda_1}^{(1)} \otimes b_{\lambda_2}^{(2)} \otimes \cdots \otimes b_{\lambda_n}^{(n)})\) , where \((\lambda_1, \ldots, \lambda_n)\) runs through \(\prod_{1 \leqslant i \leqslant n} \mathbf{L}_i\) , is a basis of \(\bigotimes_{1 \leqslant i \leqslant n} \mathbf{E}_i\) , sometimes called the tensor product of the bases \((b_{\lambda_1}^{(1)})\) in question.

Remarks. (1) The above remarks concerning the case of modules over a commutative ring generalize as in no. 5, Remark 2 when there is a tensor product \(\mathbf{E}_1 \otimes_{\mathbf{A}_1} \mathbf{E}_2 \otimes \cdots \otimes_{\mathbf{A}_{n-1}} \mathbf{E}_n\) where the rings \(\mathbf{A}_i\) are not necessarily commutative and where, for each \(i\) , there is a homomorphism \(\rho_i: \mathbf{C} \to \mathbf{A}_i\) of the same commutative ring \(\mathbf{C}\) such that: (1) \(\rho_i(\mathbf{C})\) is contained in the centre of \(\mathbf{A}_i\) ; (2) for \(2 \leqslant i \leqslant n - 1\) , the C-module structures on \(\mathbf{E}_i\) obtained using the homomorphisms \(\rho_{i-1}\) and \(\rho_i\) coincide. Then a C-module structure is obtained on \(\mathbf{E}_1 \otimes_{\mathbf{A}_1} \mathbf{E}_2 \otimes \cdots \otimes_{\mathbf{A}_{n-1}} \mathbf{E}_n\) and canonical mappings analogous to (13) (no. 5), which will be left to the reader to describe.

\begin{enumerate}
\def\labelenumi{(\arabic{enumi})}
\setcounter{enumi}{1}
\tightlist
\item
  Let A, B be two rings, E a right A-module, \(E'\) a left A-module, F a right B-module and \(F'\) a left B-module. The Z-bilinear mappings of \((\mathbf{E} \otimes_{\mathbf{A}} \mathbf{E}') \times (\mathbf{F} \otimes_{\mathbf{B}} \mathbf{F}')\) into a Z-module G are then in one-to-one correspondence with the Z-multilinear mappings f of \(E \times E' \times F \times F'\) into G satisfying the conditions
\end{enumerate}

\[
\left\{ \begin{array}{l} f (x \lambda , x ^ {\prime}, y, y ^ {\prime}) = f (x, \lambda x ^ {\prime}, y, y ^ {\prime}) \\ f (x, x ^ {\prime}, y \mu , y ^ {\prime}) = f (x, x ^ {\prime}, y, \mu y ^ {\prime}) \end{array} \right.\tag{38}
\]

for \(\lambda\in\mathbf{A}\) , \(\mu\in\mathbf{B}\) , \(x\in\mathbf{E}\) , \(x'\in\mathbf{E}'\) , \(y\in\mathbf{F}\) , \(y'\in\mathbf{F}'\) . The general constructions given in this no. reduce the proof of this to defining a canonical \(\mathbf{Z}\) -module isomorphism between \((\mathbf{E}\otimes_{\mathbf{A}}\mathbf{E}')\otimes_{\mathbf{Z}}(\mathbf{F}\otimes_{\mathbf{B}}\mathbf{F}')\) and \(\mathbf{E}\otimes_{\mathbf{A}}\mathbf{E}'\otimes_{\mathbf{Z}}\mathbf{F}\otimes_{\mathbf{B}}\mathbf{F}'\) , which follows from the associativity property of tensor products of the form (33).

\section{RELATIONS BETWEEN TENSOR PRODUCTS AND HOMOMORPHISM MODULES}

\subsection{\texorpdfstring{THE ISOMORPHISMS \(\operatorname{Hom}_{\mathbf{B}}(\mathbf{E} \otimes_{\mathbf{A}} \mathbf{F}, \mathbf{G}) \to \operatorname{Hom}_{\mathbf{A}}(\mathbf{F}, \operatorname{Hom}_{\mathbf{B}}(\mathbf{E}, \mathbf{G}))\) AND \(\operatorname{Hom}_{\mathbf{C}}(\mathbf{E} \otimes_{\mathbf{A}} \mathbf{F}, \mathbf{G}) \to \operatorname{Hom}_{\mathbf{A}}(\mathbf{E}, \operatorname{Hom}_{\mathbf{C}}(\mathbf{F}, \mathbf{G}))\)}{THE ISOMORPHISMS \textbackslash operatorname\{Hom\}\_\{\textbackslash mathbf\{B\}\}(\textbackslash mathbf\{E\} \textbackslash otimes\_\{\textbackslash mathbf\{A\}\} \textbackslash mathbf\{F\}, \textbackslash mathbf\{G\}) \textbackslash to \textbackslash operatorname\{Hom\}\_\{\textbackslash mathbf\{A\}\}(\textbackslash mathbf\{F\}, \textbackslash operatorname\{Hom\}\_\{\textbackslash mathbf\{B\}\}(\textbackslash mathbf\{E\}, \textbackslash mathbf\{G\})) AND \textbackslash operatorname\{Hom\}\_\{\textbackslash mathbf\{C\}\}(\textbackslash mathbf\{E\} \textbackslash otimes\_\{\textbackslash mathbf\{A\}\} \textbackslash mathbf\{F\}, \textbackslash mathbf\{G\}) \textbackslash to \textbackslash operatorname\{Hom\}\_\{\textbackslash mathbf\{A\}\}(\textbackslash mathbf\{E\}, \textbackslash operatorname\{Hom\}\_\{\textbackslash mathbf\{C\}\}(\textbackslash mathbf\{F\}, \textbackslash mathbf\{G\}))}}

Let E be a right A-module, F a left A-module, G a Z-module and H the Z-module of mappings \(f: E \times F \rightarrow G\) which are Z-bilinear and satisfy

\[
f (x \lambda , y) = f (x, \lambda y) \quad \text { for } x \in \mathrm{E}, y \in \mathrm{F}, \lambda \in \mathrm{A}.\tag{1}
\]

It has been seen (§ 3, no. 1, Proposition 1) that there exists a canonical Z-module homomorphism

\[
\mathrm{H} \rightarrow \operatorname{Hom} _ {\mathbf {Z}} (\mathrm{E} \otimes_ {\mathrm{A}} \mathrm{F}, \mathrm{G}).\tag{2}
\]

On the other hand, a left A-module structure has been defined on \(\mathrm{Hom}_{\mathbf{Z}}(\mathrm{E},\mathrm{G})\) and a right A-module structure on \(\mathrm{Hom}_{\mathbf{Z}}(\mathrm{F},\mathrm{G})\) (§ 3, no. 3); we may therefore consider the Z-modules \(\mathrm{Hom}_{\mathbf{A}}(\mathrm{E},\mathrm{Hom}_{\mathbf{Z}}(\mathrm{F},\mathrm{G}))\) and \(\mathrm{Hom}_{\mathbf{A}}(\mathrm{F},\mathrm{Hom}_{\mathbf{Z}}(\mathrm{E},\mathrm{G}))\) . A mapping \(f\) of \(\mathbf{E} \times \mathbf{F}\) into \(\mathbf{G}\) is canonically identified with a mapping of \(\mathbf{E}\) into the set \(\mathbf{G}^{\mathbf{F}}\) of mappings of \(\mathbf{F}\) into \(\mathbf{G}\) (Set Theory, II, § 5, no. 2); by expressing the fact that the latter mapping belongs to \(\mathrm{Hom}_{\mathbf{A}}(\mathrm{E},\mathrm{Hom}_{\mathbf{Z}}(\mathrm{F},\mathrm{G}))\) , we obtain precisely the fact that \(f\) is biadditive and conditions (1); whence there is a canonical isomorphism

\[
\mathrm{H} \rightarrow \operatorname{Hom} _ {\mathbf {A}} (\mathrm{E}, \operatorname{Hom} _ {\mathbf {Z}} (\mathrm{F}, \mathrm{G}))\tag{3}
\]

and similarly there is defined a canonical isomorphism

\[
\mathrm{H} \rightarrow \operatorname{Hom} _ {\mathbf {A}} (\mathrm{F}, \operatorname{Hom} _ {\mathbf {Z}} (\mathrm{E}, \mathrm{G})).\tag{4}
\]

Suppose now that E and G also have left (resp. right) B-module structures and that the A-module and B-module structures on E are compatible. Then \(E \otimes_{A} F\) has canonically a left (resp. right) B-module structure ( \(\S 3\) , no. 4) and on the other hand \(\operatorname{Hom}_{\mathbf{B}}(\mathbf{E}, \mathbf{G})\) has canonically a left A-module structure ( \(\S 1\) , no. 14). We may therefore consider the Z-modules \(\operatorname{Hom}_{\mathbf{B}}(\mathbf{E} \otimes_{\mathbf{A}} \mathbf{F}, \mathbf{G})\) and \(\operatorname{Hom}_{\mathbf{A}}(\mathbf{F}, \operatorname{Hom}_{\mathbf{B}}(\mathbf{E}, \mathbf{G}))\) , which are submodules of \(\operatorname{Hom}_{\mathbf{Z}}(\mathbf{E} \otimes_{\mathbf{A}} \mathbf{F}, \mathbf{G})\) and \(\operatorname{Hom}_{\mathbf{A}}(\mathbf{F}, \operatorname{Hom}_{\mathbf{Z}}(\mathbf{E}, \mathbf{G}))\) respectively ( \(\S 2\) , no. 1, Theorem 2). We examine under what condition a mapping \(f \in H\) has as image under the isomorphisms (2) and (4) an element of \(\operatorname{Hom}_{\mathbf{B}}(\mathbf{E} \otimes_{\mathbf{A}} \mathbf{F}, \mathbf{G})\) and an element of \(\operatorname{Hom}_{\mathbf{A}}(\mathbf{E}, \operatorname{Hom}_{\mathbf{B}}(\mathbf{F}, \mathbf{G}))\) respectively; in each of the two cases we find the same condition

\[
\begin{array}{r l} f (\beta x, y) & = \beta f (x, y) \\ (\text { resp.   } f (x \beta , y) & = f (x, y) \beta) \end{array}
\]

for \(x \in \mathbf{E}, y \in \mathbf{F}, \beta \in \mathbf{B}\) .

Similarly, suppose that F and G are left (resp. right) C-modules and that the A-module and C-module structures on F are compatible. Then, for a mapping \(f \in H\) to have as image under (2) or (3) an element of \(\operatorname{Hom}_{\mathbf{C}}(\mathbf{E} \otimes_{\mathbf{A}} \mathbf{F}, \mathbf{G})\) or \(\operatorname{Hom}_{\mathbf{A}}(\mathbf{E}, \operatorname{Hom}_{\mathbf{C}}(\mathbf{F}, \mathbf{G}))\) respectively, it is necessary and sufficient that it satisfy the same condition

\[
\begin{array}{c} f (x, \gamma y) = \gamma f (x, y) \\ (\text { resp. } f (x, y \gamma) = f (x, y) \gamma) \end{array}
\]

for \(x \in \mathbf{E}, y \in \mathbf{F}, \gamma \in \mathbf{C}\) .

We have therefore established the following result (in the notation introduced above):

PROPOSITION 1. (a) Let E be a (B, A)-bimodule, F a left A-module and G a left B-module. For every mapping \(g \in \operatorname{Hom}_{\mathbf{B}}(\mathbf{E} \otimes_{\mathbf{A}} \mathbf{F}, \mathbf{G})\) , let \(g'\) be the mapping of F into \(\operatorname{Hom}_{\mathbf{B}}(\mathbf{E}, \mathbf{G})\) defined by \((g'(y))(x) = g(x \otimes y)\) for \(x \in \mathbf{E}, y \in \mathbf{F}\) . The mapping \(g \mapsto g'\) is an isomorphism

\[
\beta : \operatorname{Hom} _ {\mathbf {B}} (\mathbf {E} \otimes_ {\mathbf {A}} \mathbf {F}, \mathbf {G}) \rightarrow \operatorname{Hom} _ {\mathbf {A}} (\mathbf {F}, \operatorname{Hom} _ {\mathbf {B}} (\mathbf {E}, \mathbf {G})).\tag{5}
\]

\begin{enumerate}
\def\labelenumi{(\alph{enumi})}
\setcounter{enumi}{1}
\tightlist
\item
  Let \(\mathbf{E}\) be a right \(\mathbf{A}\) -module, \(\mathbf{F}\) an \((\mathbf{A}, \mathbf{C})\) -bimodule and \(\mathbf{G}\) a right \(\mathbf{C}\) -module. For every mapping \(h \in \operatorname{Hom}_{\mathbf{C}}(\mathbf{E} \otimes_{\mathbf{A}} \mathbf{F}, \mathbf{G})\) , let \(h'\) be the mapping of \(\mathbf{E}\) into \(\operatorname{Hom}_{\mathbf{C}}(\mathbf{F}, \mathbf{G})\) defined by \((h'(x))(y) = h(x \otimes y)\) for \(x \in \mathbf{E}\) , \(y \in \mathbf{F}\) . The mapping \(h \mapsto h'\) is an isomorphism
\end{enumerate}

\[
\gamma : \operatorname{Hom} _ {\mathbf {C}} (\mathrm{E} \otimes_ {\mathbf {A}} \mathrm{F}, \mathrm{G}) \rightarrow \operatorname{Hom} _ {\mathbf {A}} (\mathrm{E}, \operatorname{Hom} _ {\mathbf {C}} (\mathrm{F}, \mathrm{G})).\tag{6}
\]

In particular B and C may be taken to be a subring \(\Gamma\) of the centre of the ring A; then for every \(\Gamma\) -module G, the three \(\Gamma\) -modules

\[
\operatorname{Hom} _ {\Gamma} (\mathrm{E} \otimes_ {\mathrm{A}} \mathrm{F}, \mathrm{G}), \quad \operatorname{Hom} _ {\mathrm{A}} (\mathrm{E}, \operatorname{Hom} _ {\Gamma} (\mathrm{F}, \mathrm{G})), \quad \operatorname{Hom} _ {\mathrm{A}} (\mathrm{F}, \operatorname{Hom} _ {\Gamma} (\mathrm{E}, \mathrm{G}))
\]

are canonically isomorphic to the \(\Gamma\) -module of \(\Gamma\) -bilinear mappings of \(\mathbf{E} \times \mathbf{F}\) into \(\mathbf{G}\) which satisfy (1). More particularly:

COROLLARY. If C is a commutative ring and E, F, G three C-modules, then the C-modules

\[
\begin{array}{l l} \operatorname{Hom} _ {\mathbf {c}} (\mathrm{E} \otimes_ {\mathbf {c}} \mathrm{F}, \mathrm{G}), & \operatorname{Hom} _ {\mathbf {c}} (\mathrm{E}, \operatorname{Hom} _ {\mathbf {c}} (\mathrm{F}, \mathrm{G})), \\ \operatorname{Hom} _ {\mathbf {c}} (\mathrm{F}, \operatorname{Hom} _ {\mathbf {c}} (\mathrm{E}, \mathrm{G})), & \mathcal {L} _ {2} (\mathrm{E}, \mathrm{F}; \mathrm{G}) \end{array}
\]

are canonically isomorphic.
